%% ============================================================
%% Oblate Spheroid Excitation Theory (OSET): Acta Materialia
%% Two-column (elsarticle 5p), full derivations retained
%% pdflatex -> bibtex -> pdflatex -> pdflatex
%% ============================================================
\documentclass[5p,sort&compress]{elsarticle}

\usepackage{amsmath,amssymb,amsfonts}
\usepackage{mathtools}
\usepackage{graphicx}
\usepackage{booktabs}
\usepackage{array}
\usepackage{xcolor}
\usepackage{siunitx}
\usepackage[hidelinks]{hyperref}
\usepackage{microtype}

\graphicspath{{figures/}}
\journal{Acta Materialia}

\newcommand{\OSTZ}{\textsc{ostz}}
\newcommand{\OSET}{\textsc{oset}}
\newcommand{\dF}{\Delta F_0}
\newcommand{\gz}{\gamma_0}
\newcommand{\beff}{b_\mathrm{eff}}
\newcommand{\PN}{Peierls--Nabarro}
\newcommand{\half}{\tfrac{1}{2}}

\begin{document}
\begin{frontmatter}

\title{Oblate Spheroid Excitation Theory (\OSET): A Unified,
  Lattice-Free Foundation for Plastic Deformation from Which
  Dislocations Emerge as Collective Excitations}

\author[iitk]{Albert~Linda\corref{cor}}
\ead{albert.hzbn@gmail.com}
\author[anna]{K.A.~Padmanabhan\corref{cor}}
\ead{ananthaster@gmail.com}

\cortext[cor]{Corresponding author.}
\affiliation[iitk]{organization={Department of Materials Science and
  Engineering, Indian Institute of Technology Kanpur},
  city={Kanpur}, postcode={208016}, country={India}}
\affiliation[anna]{organization={Professor of Eminence, Materials
  Science and Engineering Programme, Department of Mechanical
  Engineering, Anna University},
  city={Chennai}, postcode={600\,025}, country={India}}

\begin{abstract}
Dislocation theory has underpinned crystal plasticity for nearly a
century, yet its central object, a topological line defect defined
through a Burgers circuit in a Bravais lattice, cannot describe
plastic flow in grain boundaries, metallic glasses, ceramics,
polymers, or nanocrystals approaching the glass transition, simply
because no periodic lattice is available there. We propose the
\emph{Oblate Spheroid Excitation Theory} (\OSET), in which the
elementary carrier of plastic deformation in any solid, crystalline
or not, is a shear-eigenstrained oblate spheroid, the
oblate-spheroidal transformation zone (\OSTZ), treated exactly within
Eshelby's inclusion theory. The \OSTZ\ requires no lattice and has a
finite, non-singular activation energy and stress field that is
intrinsically thermally activated rather than appended afterwards.
Three results are proved: a single \OSTZ\ produces a non-singular
elastic dipole rather than a dislocation's singular field; a
co-planar chain of $N$ \OSTZ s is mathematically identical to a
Peierls--Nabarro dislocation, with both its core width and its
Burgers vector fixed by \OSTZ\ geometry; and a genuine lattice
dislocation nucleates only once the chain reaches a critical length
set by the host lattice. Dislocations therefore emerge as a
collective, large-$N$ limit of \OSET\ rather than being assumed as an
independent entity, and the theoretical shear strength, Peierls
stress, dislocation core energy, Frank--Read critical stress, and
stacking-fault energy all follow as derived, parameter-free
quantities. \OSET\ is validated against grain-boundary-sliding data, independent literature spanning
metals, ceramics, and bulk metallic glasses, and a recent 41-system
compilation, reproducing the fitted dilatational eigenstrain to
within 2\% and the shear eigenstrain to within 15\% with no
adjustable parameters. Because classical dislocation theory emerges
from \OSET\ but \OSET\ does not require dislocations, \OSET\
provides a more fundamental and more broadly applicable foundation
for plastic deformation across material classes.
\end{abstract}

\begin{keyword}
oblate spheroid excitation \sep Eshelby inclusion \sep dislocation
emergence \sep grain-boundary sliding \sep Peierls--Nabarro model \sep
superplasticity
\end{keyword}
\end{frontmatter}

%%=========================================================
\section{Introduction}
\label{sec:intro}

The dislocation is, by a wide margin, the most successful concept in
the theory of crystal plasticity. Since its independent introduction
by Taylor, Orowan, and Polanyi in 1934~\cite{taylor1934,orowan1934,polanyi1934}
to explain why real crystals yield at stresses orders of magnitude
below the theoretical shear strength, the dislocation has provided
the microscopic basis for strain hardening, the Hall--Petch grain-size
dependence of strength~\cite{hall1951,petch1953}, creep, fatigue, and
fracture in essentially every crystalline engineering alloy. Its
mathematical foundation was laid by Volterra's elastic continuum
construction of a line defect~\cite{volterra1907} decades before its
physical role was understood, and was subsequently refined by the
Peierls--Nabarro model~\cite{peierls1940,nabarro1947}, which regularised
the otherwise singular core by spreading the displacement
discontinuity over a finite width, and by Eshelby's general theory of
elastic inclusions~\cite{eshelby1957}, which placed the energetics of
lattice defects on a rigorous continuum-mechanical footing. Eight
decades of refinement have made dislocation theory, augmented with
kink-pair nucleation, climb, cross-slip, and forest-hardening models,
the standard language in which crystal plasticity is taught,
computed, and engineered~\cite{hirth1982}.

That success rests, however, on a definition that is fundamentally
topological rather than physical. The dislocation is specified by
the closure failure of a Burgers circuit traversed in a Bravais
lattice,
\begin{equation}
  \oint_C \mathrm{d}u_i = b_i,\qquad b_i\in\Lambda_\mathrm{Bravais},
  \label{eq:burgers}
\end{equation}
and this definition carries three restrictions that are unavoidable
consequences of its topological character rather than incidental
simplifications. \textbf{(R1)} It requires a crystal lattice: without
a Bravais lattice there is no notion of a closed reference circuit,
no quantised Burgers vector $\Lambda_\mathrm{Bravais}$, and therefore
no dislocation, so grain boundaries, glasses, and polymers are
excluded from the description by definition rather than by
approximation. \textbf{(R2)} Its core is singular: the Volterra
stress field $\sigma\sim Gb/2\pi r$ diverges as $r\to0$, forcing the
introduction of an empirical cutoff radius $r_0\sim b$ and a core
energy that must be fitted to experiment or atomistic simulation
rather than derived. \textbf{(R3)} It is athermal at the elementary
level: the bare dislocation carries no intrinsic activation energy or
entropy, so thermal activation has to be appended afterwards through
kink-nucleation or Peierls--Nabarro overlays rather than emerging from
the definition itself.

These three restrictions are not merely formal concerns; each one
fails in a regime of substantial current interest. In grain-boundary
sliding and superplastic deformation, the dislocation content of a
boundary is conventionally described through DSC
(displacement-shift-complete) lattice vectors, but for a general
high-angle boundary the DSC lattice is commensurate with the two
adjoining crystals only in a statistical sense, so the underlying
reference lattice that R1 demands is simply not available; this is
precisely the regime in which Padmanabhan and
co-workers~\cite{pad-part1,padgleiter2014,paddivinski2024} introduced a
phenomenological rate equation, built around an empirical
boundary-sliding eigenstrain, that successfully predicts superplastic
flow without ever invoking a dislocation explicitly. In nanocrystalline
metals with grain size $d<10b$, standard dislocation-based hardening
theory predicts a continued increase in strength with decreasing grain
size that is not observed experimentally; instead, below a critical
size the flow stress falls with further grain refinement, the inverse
Hall--Petch effect~\cite{schiotz1998}, because a grain too small to
host an extended dislocation network cannot be described by a theory
whose elementary unit is a dislocation line. In metallic glasses there
is no lattice at all, and yet the material deforms plastically; the
accepted microscopic picture is the shear-transformation zone
(STZ)~\cite{argon1979,falk-langer1998}, a local, thermally activated
rearrangement of a cluster of atoms that bears no resemblance to a
topological line defect. Ceramics and semicrystalline polymers present
further intermediate cases: wherever dislocation glide is restricted
to a small number of slip systems or absent altogether, the same
absence of a usable reference lattice (R1) again forces boundary- and
free-volume-mediated mechanisms to dominate the observed plasticity.

The existing remedies for these failures are mechanism-specific rather
than unified. STZ theory~\cite{argon1979,falk-langer1998} supplies a
phenomenological activation volume and energy for amorphous flow but
does not derive either from an explicit elastic inclusion, so it
offers no route back to dislocation quantities such as the Peierls
stress or stacking-fault energy. The Padmanabhan et al.\ grain-boundary
sliding theory~\cite{pad-part1} is, in the present context, the
closest predecessor to this work: it identifies the elementary
boundary-sliding event with an oblate-spheroidal region and obtains a
rate equation in excellent agreement with experiment, but its
activation energy is expressed in terms of Eshelby constraint factors
that were asserted rather than derived from the underlying $I$-integrals,
and the theory was never connected back to the dislocation quantities
of single-crystal plasticity. What is missing, in short, is a single
elementary excitation that (i) requires no crystal lattice for its
definition, (ii) has a finite, non-singular core energy obtained from
elasticity theory rather than fitted to data, (iii) is intrinsically
thermally activated, and (iv) reduces to the conventional dislocation,
with all of its derived properties, in the appropriate collective
limit.

This paper proposes such an excitation: the \emph{oblate spheroidal
transformation zone} (\OSTZ), and develops it into a complete
framework that we call the \emph{Oblate Spheroid Excitation Theory}
(\OSET). Physically, an \OSTZ\ is the smallest unit of irreversible
shear a solid can sustain without invoking a lattice: a disk-shaped
region, of the same size as a free-volume cluster at a grain boundary
or a shear-transformation zone in a glass, that flips into a sheared
configuration when thermally or mechanically activated. Modelling
this region as an Eshelby oblate spheroid, rather than as a sphere or
a flat crack, is the key geometric choice of the theory: the shape is
thin enough to localise shear near a glide plane, yet smooth enough
that Eshelby's inclusion theory~\cite{eshelby1957,mura1987} yields an
exact, closed-form, and non-singular stress and energy. \OSET\
requires exactly three material inputs, the \OSTZ\ radius $W$, the
shear eigenstrain $\gz$, and the shear modulus $G$ of the host, and
from these derives every other quantity used in this paper.

Three central results are established below. First, the far field of
a single \OSTZ\ is shown to be a stress \emph{dipole} rather than a
monopole (\S\ref{sec:dipole}), decaying as $r^{-3}$ and therefore
incapable, on its own, of producing the long-range $r^{-1}$ field of a
dislocation. Second, a co-planar chain of $N$ \OSTZ s is proved to be
mathematically identical to a Peierls--Nabarro dislocation
(\S\ref{sec:emergence}), with core half-width $\zeta=W$ and Burgers
vector $b=N\gz W$ following directly from the chain construction
rather than being assumed. Third, a full lattice dislocation is shown
to nucleate only once the chain reaches a critical length
$N=N_c=b_\mathrm{latt}/(\gz W)$, at which point the Burgers vector
locks onto a Bravais lattice vector. Because the dislocation emerges
from \OSTZ\ mechanics, rather than the reverse, \OSET\ further yields,
as derived rather than fitted quantities, the theoretical shear
strength, the Peierls stress, the dislocation core energy, the
Frank--Read critical stress, and the stacking-fault energy
(\S\ref{sec:derived}), together with a grain-boundary energy
expression (\S\ref{sec:gb}) and a statistical-mechanical treatment of
the OSTZ population that recovers the Padmanabhan et al.\ superplastic
rate equation and a Taylor-hardening law (\S\ref{sec:stat}). Sections
\ref{sec:classes}--\ref{sec:validation} extend the framework across
material classes, state its novel testable predictions, and validate
it quantitatively against the original 1996 Padmanabhan et al.\ data
set, subsequent independent literature, and the 41-system compilation
of Harisankar and Padmanabhan~\cite{harisankar2025}.

%%=========================================================
\section{The Fundamental Entity: the \OSTZ}
\label{sec:ostz}

Physically, an \OSTZ\ is the smallest unit of irreversible shear
that a solid can sustain without invoking a crystal lattice: a
roughly disk-shaped pocket of material, of the same size as a
free-volume cluster at a grain boundary or a shear-transformation
zone in a glass, that flips into a sheared configuration when
thermally or mechanically activated. Modelling it as an
\emph{oblate spheroid} rather than a sphere or a flat crack is the
key geometric choice of this paper: it is thin enough to represent a
localized shear event confined near a boundary or free-volume site,
yet smooth enough that Eshelby's inclusion theory~\cite{eshelby1957}
gives an exact, closed-form, non-singular stress and energy. The
next two subsections fix this geometry and then derive its
energetics from first principles.

\subsection{Geometric definition}
\label{sec:geometry}

Only two numbers are needed to fix the shape of the \OSTZ: its
in-plane radius $W$, set by the free-volume length scale of the
host material, and its aspect ratio $\alpha$, fixed once and for all
by requiring the spheroid to be oblate enough to localise shear near
a glide plane while remaining smooth enough for Eshelby's exact
solution to apply. We adopt $\alpha=1/2$ throughout: an \OSTZ\ is an
oblate spheroid that has undergone a shear eigenstrain, with
semi-axes
\begin{align}
  a_1=a_2=W,\quad a_3=W/2,&\quad \alpha\equiv a_3/a_1=\tfrac{1}{2},
  \notag\\
  V_0=\tfrac{4\pi}{3}a_1^2 a_3&=\tfrac{2\pi}{3}W^3.
  \label{eq:geometry}
\end{align}
The flat face lies in the $x_1$--$x_2$ glide plane, so that shear
parallel to this plane corresponds to relative sliding of the two
flat faces of the spheroid past each other, the same kinematics
as grain-boundary sliding or a shear-transformation event. The
eigenstrain that produces this sliding, together with the
accompanying volume change measured by $\varepsilon_0$, is
\begin{equation}
  \varepsilon^*_{13}=\varepsilon^*_{31}=\gz/2\;(\text{shear}),\quad
  \varepsilon^*_{ii}=\varepsilon_0/3\;(\text{dilat.})
  \label{eq:eigenstrain}
\end{equation}

Because the eigenstrain is a pure $13$-shear rather than an in-plane
($12$) shear, this geometry makes $S_{1313}$, not $S_{1212}$, the
Eshelby tensor component governing the stored elastic energy:
the central quantity computed in the rest of this section.

\subsection{Eshelby solution for the \OSTZ}
\label{sec:eshelby}

We now need the actual stress and strain that the eigenstrain of
Eq.~\eqref{eq:eigenstrain} produces inside the \OSTZ, since this is
what determines how much elastic energy is stored when the \OSTZ\
activates (\S\ref{sec:energy}). For an arbitrary inclusion shape this
would require solving an elasticity boundary-value problem
numerically; the entire point of choosing an \emph{ellipsoid} is that
Eshelby's inclusion theorem~\cite{eshelby1957} gives the answer in
closed form. The theorem states that the strain and stress inside a
homogeneous ellipsoidal inclusion are \emph{uniform} (independent of
position within the inclusion):
\begin{align}
  \varepsilon_{ij}^\mathrm{in}&=S_{ijkl}\varepsilon^*_{kl},
    \label{eq:eshelby1}\\
  \sigma_{ij}^\mathrm{in}&=C_{ijkl}(S_{klmn}-I_{klmn})\varepsilon^*_{mn}.
    \label{eq:eshelby2}
\end{align}
All of the shape-dependence is packaged into one object, the Eshelby
tensor $S_{ijkl}$; uniformity of $\varepsilon^\mathrm{in}_{ij}$ holds
only for ellipsoids, because they are the only shapes for which
$S_{ijkl}$ does not vary from point to point inside the inclusion.

Since the \OSTZ\ eigenstrain (Eq.~\eqref{eq:eigenstrain}) is a pure
$13$-shear, only the single component $S_{1313}$ is needed; for the
general ellipsoid this component is built from two auxiliary shape
integrals, $I_\alpha$ and $I_{\alpha\beta}$, that depend only on the
semi-axes~\cite{mura1987}:
\begin{equation}
  S_{1313}=\frac{(1+\alpha^2)I_{13}+(1-2\nu)(I_1+I_3)}{16\pi(1-\nu)},
  \label{eq:S1313}
\end{equation}
where $\Delta(s)=\sqrt{(a_1^2{+}s)(a_2^2{+}s)(a_3^2{+}s)}$,
\begin{align}
  I_\alpha&=2\pi a_1 a_2 a_3\int_0^\infty
    \frac{\mathrm ds}{(a_\alpha^2+s)\Delta(s)},
    \label{eq:Ialpha}\\
  I_{\alpha\beta}&=2\pi a_1 a_2 a_3\int_0^\infty
    \frac{\mathrm ds}{(a_\alpha^2+s)(a_\beta^2+s)\Delta(s)},\notag
\end{align}
with sum rule $I_1+I_2+I_3=4\pi$. These integrals are elementary but
tedious; the next three subsections evaluate them explicitly for our
oblate spheroid ($a_1=a_2=1$, $a_3=\alpha$ in reduced units), first
deriving $I_1$ in closed form, then $I_{13}$, then substituting
numbers for $\alpha=1/2$ to obtain the constraint factor $\beta_1$
that ultimately enters the activation energy of Eq.~\eqref{eq:dF0}.

\subsubsection{Derivation of $I_1$}
\label{sec:I1}

\paragraph{Starting point}
Specialize to $a_1=a_2=1$, $a_3=\alpha$. Then
$\Delta(s)=(1+s)\sqrt{\alpha^2+s}$ and
\begin{equation}
  I_1=2\pi\alpha\int_0^\infty\frac{\mathrm ds}
  {(1+s)^2\sqrt{\alpha^2+s}}.
  \label{eq:I1start}
\end{equation}

\paragraph{Substitution}
Let $t=\sqrt{\alpha^2+s}$, so $s=t^2-\alpha^2$, $\mathrm ds=2t\,\mathrm dt$,
$\sqrt{\alpha^2+s}=t$, and $1+s=t^2+(1-\alpha^2)\equiv t^2+e^2$
(where $e^2\equiv1-\alpha^2$). New limits: $t:\alpha\to\infty$.
Substituting:
\begin{equation}
  I_1=4\pi\alpha\int_\alpha^\infty\frac{\mathrm dt}{(t^2+e^2)^2}.
  \label{eq:I1sub}
\end{equation}

\paragraph{Antiderivative}
The reduction formula (proved by differentiating the right side):
\begin{equation}
  \int\frac{\mathrm dt}{(t^2+e^2)^2}=
  \frac{t}{2e^2(t^2+e^2)}+\frac{\arctan(t/e)}{2e^3}+C.
  \label{eq:antideriv}
\end{equation}
\paragraph{Proof} Let $F=t/[2e^2(t^2{+}e^2)]$,
$G_f=\arctan(t/e)/(2e^3)$:
\begin{align*}
  F'&=\frac{1}{2e^2}\cdot\frac{(t^2+e^2)-t\cdot2t}{(t^2+e^2)^2}
    =\frac{e^2-t^2}{2e^2(t^2+e^2)^2},\\
  G_f'&=\frac{1}{2e^3}\cdot\frac{1/e}{1+(t/e)^2}
    =\frac{1}{2e^2(t^2+e^2)}.
\end{align*}
Sum: $F'+G_f'=[(e^2-t^2)+(t^2+e^2)]/[2e^2(t^2+e^2)^2]
=1/(t^2+e^2)^2$. $\square$

\paragraph{Evaluating at the bounds}
\textit{Upper bound} ($t\to\infty$): $F\to0$, $\arctan\to\pi/2$, so
$[F+G_f]_{t\to\infty}=\pi/(4e^3)$.

\textit{Lower bound} ($t=\alpha$): note $\alpha^2+e^2=1$, so
$F|_\alpha=\alpha/(2e^2)$,
$G_f|_\alpha=\arctan(\alpha/e)/(2e^3)$. Hence:
\begin{align}
  \int_\alpha^\infty\frac{\mathrm dt}{(t^2+e^2)^2}
  &=\frac{\pi}{4e^3}-\frac{\alpha}{2e^2}
    -\frac{\arctan(\alpha/e)}{2e^3}\notag\\
  &=\frac{1}{2e^3}\!\left[\frac{\pi}{2}-\arctan\!\left(\frac{\alpha}{e}
    \right)\right]-\frac{\alpha}{2e^2}.
  \label{eq:I1bounds}
\end{align}

\paragraph{Identity $\pi/2-\arctan(\alpha/e)=\arccos\alpha$}
Set $\alpha=\cos\phi$ ($e=\sin\phi$). Then
$\arctan(\cos\phi/\sin\phi)=\arctan(\cot\phi)=\pi/2-\phi$, so
$\pi/2-\arctan(\alpha/e)=\phi=\arccos\alpha$. Substituting into
Eq.~\eqref{eq:I1bounds} and multiplying by $4\pi\alpha$, and
replacing $e^3=(1-\alpha^2)^{3/2}$:
\begin{equation}
  \boxed{\begin{gathered}
  I_1=\frac{2\pi\alpha}{(1-\alpha^2)^{3/2}}
  \!\left[\arccos\alpha-\alpha\sqrt{1-\alpha^2}\right],\\
  I_3=4\pi-2I_1.
  \end{gathered}}
  \label{eq:I1result}
\end{equation}

\subsubsection{Derivation of $I_{13}=(I_3-I_1)/(1-\alpha^2)$}
\label{sec:I13}

Equation~\eqref{eq:S1313} for $S_{1313}$ still needs the
cross-integral $I_{13}$, which superficially requires evaluating yet
another integral of the same type as $I_1$. In fact $I_{13}$ can be
obtained algebraically from $I_1$ and $I_3$ without any new
integration, by exploiting the way the two denominators
$(1+s)$ and $(\alpha^2+s)$ combine. Writing the integrals explicitly
($\Delta=(1+s)\sqrt{\alpha^2+s}$):
\begin{align}
  I_1&=2\pi\alpha\int_0^\infty\frac{\mathrm ds}
    {(1+s)^2\sqrt{\alpha^2+s}},\notag\\
  I_3&=2\pi\alpha\int_0^\infty\frac{\mathrm ds}
    {(1+s)(\alpha^2+s)^{3/2}},\notag\\
  I_{13}&=2\pi\alpha\int_0^\infty\frac{\mathrm ds}
    {(1+s)^2(\alpha^2+s)^{3/2}}.\notag
\end{align}
Forming $I_3-I_1$ and factoring out $1/[(1+s)\sqrt{\alpha^2+s}]$:
\begin{align*}
  I_3-I_1&=2\pi\alpha\int_0^\infty\frac{1}{(1+s)\sqrt{\alpha^2+s}}
  \!\left[\frac{1}{\alpha^2+s}-\frac{1}{1+s}\right]\mathrm ds.
\end{align*}
The bracket equals $(1-\alpha^2)/[(\alpha^2{+}s)(1{+}s)]$:
\begin{align*}
  I_3-I_1&=(1-\alpha^2)\cdot2\pi\alpha\int_0^\infty
  \frac{\mathrm ds}{(1+s)^2(\alpha^2+s)^{3/2}}\\
  &=(1-\alpha^2)I_{13}.
\end{align*}
Therefore:
\begin{equation}
  \boxed{I_{13}=\frac{I_3-I_1}{1-\alpha^2}.}\quad\square
  \label{eq:I13result}
\end{equation}

\subsubsection{Numerical evaluation ($\alpha=1/2$, $\nu=1/3$)}
\label{sec:numeval}

With $I_1$ and $I_{13}$ now available in closed form, we can finally
put in the actual \OSTZ\ aspect ratio $\alpha=1/2$ and a
representative Poisson ratio $\nu=1/3$ to obtain the single number,
$\beta_1$, that controls the shear-energy cost of activating an
\OSTZ\ (Eq.~\eqref{eq:Eshear} below).

\paragraph{Step~1: compute $I_1$}
$e=\sqrt{1-0.25}=0.8660$.
Bracket: $\arccos(0.5)-0.5\times0.8660=1.0472-0.4330=0.6142$.
$(0.75)^{3/2}=0.6495$.
$I_1=\pi\times0.6142/0.6495=\mathbf{2.972}$.

\paragraph{Step~2: $I_3$ and $I_{13}$}
$I_3=4\pi-2\times2.972=12.566-5.944=\mathbf{6.622}$.
$I_{13}=(6.622-2.972)/0.75=3.650/0.75=\mathbf{4.867}$.

\paragraph{Step~3: $S_{1313}$ from Eq.~\eqref{eq:S1313}}
Numerator: $1.25\times4.867+(1-2/3)\times9.594
=6.084+3.198=9.282$.
Denominator: $16\pi\times(2/3)=33.51$.
$S_{1313}=9.282/33.51=\mathbf{0.2772}$.

\paragraph{Step~4: $\beta_1$}
$\beta_1=1-2\times0.2772=\mathbf{0.4456}\approx0.446$.

\subsubsection{Limiting-case verification}
\label{sec:limits}

Before trusting $\beta_1=0.446$ as a physically reasonable
constraint factor, it is worth checking the formula against two
shapes whose answer we already know by symmetry or by elementary
physical reasoning: the sphere, where isotropy gives an independent
check, and the vanishingly thin disk, which should behave like a
stress-free crack.

\paragraph{Sphere ($\alpha\to1$)}
For a sphere all semi-axes equal: $I_1=I_2=I_3=4\pi/3$.
Direct integration:
$I_{13}^\text{sphere}=2\pi\int_0^\infty(1+s)^{-7/2}\mathrm ds
=2\pi[u^{-5/2}/(-5/2)]_1^\infty=4\pi/5$.
Substituting ($1+\alpha^2\to2$, $I_1+I_3=8\pi/3$):
\begin{equation}
  S_{1313}^\text{sphere}=\frac{4-5\nu}{15(1-\nu)}.
  \label{eq:S1313sphere}
\end{equation}
For $\nu=1/3$: $S_{1313}^\text{sphere}=7/30=0.233$,
$\beta_1^\text{sphere}=8/15\approx0.533$.

\paragraph{Thin disk ($\alpha\to0$)}
$I_3\to4\pi$, $I_{13}\to4\pi$, $S_{1313}^\text{disk}=\half$,
$\beta_1\to0$: a crack-like disk stores no shear energy.
The OSTZ aspect ratio $\alpha=1/2$ gives the correct intermediate
(Figure~S1 in the Supporting Information plots $S_{1313}(\alpha)$
across this full range for several $\nu$).

\subsubsection{Dilatational constraint factor $\beta_2$}
\label{sec:beta2}

The shear eigenstrain is not the whole story: activating an \OSTZ\
also requires a small volume change $\varepsilon_0$ (free-volume
creation), and this dilatation costs its own elastic energy,
governed by an analogous constraint factor $\beta_2$. Because a
volume change has no preferred direction, it is consistent to model
this part of the problem with the simpler, fully isotropic spherical
inclusion rather than repeating the oblate-spheroid calculation: for
a spherical inclusion with purely dilatational eigenstrain
$\varepsilon^*_{ij}=(\varepsilon_0/3)\delta_{ij}$, the sphere Eshelby
tensor satisfies:
\begin{equation}
  S_{ijkk}^\text{sphere}=\frac{1+\nu}{3(1-\nu)}\delta_{ij}.
  \label{eq:Sijkk}
\end{equation}
This follows from the known sphere values~\cite{mura1987}
$S_{1111}=(7-5\nu)/[15(1-\nu)]$,
$S_{1122}=(5\nu-1)/[15(1-\nu)]$, summing:
$S_{1111}+S_{1122}+S_{1133}=(1+\nu)/[3(1-\nu)]$. Full derivation
of $\beta_2$ follows in \S\ref{sec:Edilat}; the result:
\begin{equation}
  \boxed{\beta_2=\frac{4(1+\nu)}{9(1-\nu)}.}
  \label{eq:beta2}
\end{equation}
For $\nu=1/3$: $\beta_2=4(4/3)/[9(2/3)]=0.889$.

Two formulas for $\beta_2$ exist: (a) exact oblate formula
$\beta_2^\text{oblate}=1-S_{3333}(1-2\nu)/[2(1-\nu)]=0.816$;
(b) spherical dilatation approximation $\beta_2^\text{sphere}=0.889$~\cite{pad-part1}.
They differ by $\sim9\%$ at $\nu=1/3$;
since $\beta_2\varepsilon_0^2\ll\beta_1\gz^2$, the difference is
$<2\%$ of $\dF$ and inconsequential.

\begin{table}[htbp]
\centering
\small
\caption{Eshelby components for $\alpha=1/2$.}
\label{tab:eshelby}
\begin{tabular}{ccccc}
\toprule
$\nu$ & $S_{1313}$ & $S_{1212}$ & $S_{3333}$ & $\beta_1$\\
\midrule
0.30 & 0.2822 & 0.1769 & 0.7264 & 0.4356\\
1/3  & 0.2772 & 0.1739 & 0.7364 & 0.4456\\
0.35 & 0.2745 & 0.1723 & 0.7418 & 0.4510\\
0.40 & 0.2656 & 0.1670 & 0.7596 & 0.4688\\
\bottomrule
\end{tabular}
\end{table}

\subsection{Note on $\beta_1$: two physical regimes}

The value $\beta_1\approx0.45$ just derived is the purely elastic
constraint factor for an \OSTZ\ embedded in a matrix of its own
shear modulus $G$. The original Padmanabhan et al.\ treatment of
grain-boundary sliding, however, used an effective value
$\beta_1^\text{eff}\approx1$, more than double the elastic result.
The two are not in conflict once the different elastic environment
of a grain-boundary \OSTZ\ is accounted for, as shown below.

\textbf{(a) Eshelby (elastic) value:}
$\beta_1^\text{Esh}=1-2S_{1313}\approx0.44$--$0.49$
for $\alpha=0.5$, $\nu=0.28$--$0.44$.

\textbf{(b) Padmanabhan et al.\ effective value:} $\beta_1^\text{eff}\approx1$.

\paragraph{Physical reconciliation} OSTZs at grain boundaries are embedded
in GB material with $G_\text{GB}\approx G/2$ (GB softening~\cite{wolf1990}):
$\beta_1^\text{eff}=\beta_1^\text{Esh}\times G/G_\text{GB}
\approx0.45\times2=0.90\approx1$.

\subsection{OSTZ elastic strain energy}
\label{sec:energy}

We can now assemble the central result of this section: the
activation energy $\dF$ that a thermal fluctuation or applied stress
must supply to flip a single \OSTZ\ into its sheared configuration.
Eshelby's theorem also gives a general formula for the elastic
energy stored by an eigenstrained inclusion, as the work done by the
(now known) internal stress against the eigenstrain:
$\dF=-\half\sigma_{ij}^\mathrm{in}\varepsilon^*_{ij}V_0$. Since the
\OSTZ\ eigenstrain has two physically distinct parts (the
$13$-shear of Eq.~\eqref{eq:eigenstrain} and the dilatation
$\varepsilon_0$), we evaluate their energy contributions
separately below and add them at the end, justified by the fact
that shear and dilatation are orthogonal modes of the isotropic
stiffness tensor (no cross term). The dilatational contribution is
the smaller of the two: for the canonical values $\gz=0.1$,
$\varepsilon_0=0.05$ it accounts for roughly 25\% of the total
$\dF$, but it is not negligible and is retained throughout.

\subsubsection{Part~1: Shear contribution}
\label{sec:Eshear}

Eigenstrain $\varepsilon^*_{13}=\varepsilon^*_{31}=\gz/2$, all other
components zero.

\paragraph{Step~1 (constrained strain)}
Eq.~\eqref{eq:eshelby1} gives $\varepsilon_{13}^\mathrm{in}$ as a
contraction of the Eshelby tensor with the \emph{full} eigenstrain
tensor, so both nonzero components $\varepsilon^*_{13}$ and
$\varepsilon^*_{31}$ contribute:
$\varepsilon_{13}^\mathrm{in}=S_{1313}\varepsilon^*_{13}+S_{1331}\varepsilon^*_{31}$.
By the minor symmetry of the Eshelby tensor, $S_{1313}=S_{1331}$ (a
direct consequence of the symmetry $\varepsilon^*_{kl}=\varepsilon^*_{lk}$
of any physical strain tensor), so the two terms combine:
$\varepsilon_{13}^\mathrm{in}=S_{1313}(\gz/2+\gz/2)=S_{1313}\gz$.

\paragraph{Step~2 (elastic strain)}
The eigenstrain $\varepsilon^*_{13}$ is, by definition, a stress-free
transformation strain; it is the additional strain
$e_{13}^\mathrm{el}=\varepsilon_{13}^\mathrm{in}-\varepsilon^*_{13}$,
the difference between the actual (constrained) strain and the
eigenstrain, that the surrounding matrix elastically resists and
that therefore generates stress via Hooke's law:
$e_{13}^\mathrm{el}=\varepsilon_{13}^\mathrm{in}-\varepsilon^*_{13}
=(S_{1313}-\half)\gz$.

\paragraph{Step~3 (interior stress)}
For an isotropic matrix the relevant stiffness component is simply
twice the shear modulus, $C_{1313}=C_{1331}=G$ (the standard
isotropic-elasticity relation between shear stress and shear strain),
so
$\sigma_{13}^\mathrm{in}=2Ge_{13}^\mathrm{el}
=2G(S_{1313}-\half)\gz$.

\paragraph{Step~4 (energy contraction)}
Both $(1,3)$ and $(3,1)$ contribute equally:
\[
  \dF^\text{shear}=-\half(\sigma_{13}^\mathrm{in}\varepsilon^*_{13}
  +\sigma_{31}^\mathrm{in}\varepsilon^*_{31})V_0
  =-\half\sigma_{13}^\mathrm{in}\gz V_0.
\]
Substituting:
$\dF^\text{shear}=-\half\cdot2G(S_{1313}-\half)\gz\cdot\gz V_0
=G(\half-S_{1313})\gz^2 V_0$.
\begin{equation}
  \boxed{\dF^\text{shear}=\half\beta_1 G\gz^2 V_0,\quad
  \beta_1\equiv1-2S_{1313}.}
  \label{eq:Eshear}
\end{equation}
Since $S_{1313}<\half$ for any physical oblate spheroid,
$\beta_1>0$ and the shear energy is positive definite.

\subsubsection{Part~2: Dilatational contribution}
\label{sec:Edilat}

The same energy bookkeeping is now repeated for the volumetric part
of the eigenstrain, using the spherical-inclusion Eshelby result
$S_{ijkk}^\text{sphere}$ already obtained in
\S\ref{sec:beta2} (Eq.~\eqref{eq:Sijkk}).
Eigenstrain $\varepsilon^*_{ij}=(\varepsilon_0/3)\delta_{ij}$.

\paragraph{Step~1 (constrained strain inside a sphere)}
From Eqs.~\eqref{eq:eshelby1} and \eqref{eq:Sijkk}:
\begin{equation}
  \varepsilon_{ij}^\mathrm{in}=S_{ijkl}\varepsilon^*_{kl}
  =\frac{\varepsilon_0}{3}\cdot\frac{1+\nu}{3(1-\nu)}\delta_{ij}
  =\frac{(1+\nu)\varepsilon_0}{9(1-\nu)}\delta_{ij}.
  \label{eq:constrained_dil}
\end{equation}

\paragraph{Step~2 (elastic strain)}
\begin{align}
  e_{ij}^\mathrm{el}&=\varepsilon_{ij}^\mathrm{in}-\varepsilon^*_{ij}
    =\varepsilon_0\!\left[\frac{1+\nu}{9(1-\nu)}-\frac{1}{3}\right]\delta_{ij}
    \notag\\
  &=\varepsilon_0\cdot\frac{(1+\nu)-3(1-\nu)}{9(1-\nu)}\delta_{ij}
    =-\frac{2(1-2\nu)\varepsilon_0}{9(1-\nu)}\delta_{ij}.
    \label{eq:edilelastic}
\end{align}
(Numerator: $1+\nu-3+3\nu=4\nu-2=-2(1-2\nu)$.)
Volumetric trace:
$e_{kk}^\mathrm{el}=-6(1-2\nu)\varepsilon_0/[9(1-\nu)]
=-2(1-2\nu)\varepsilon_0/[3(1-\nu)]$.

\paragraph{Step~3 (interior stress)}
Hooke's law for an isotropic solid splits the stress into a
hydrostatic part proportional to the volumetric strain $e_{kk}^\mathrm{el}$
(governed by the Lamé constant $\lambda$) and a deviatoric part
proportional to the local strain component itself (governed by
$2G$); both parts must be added to get the total interior stress.
Using $\lambda=2G\nu/(1-2\nu)$:
\begin{align}
  \lambda e_{kk}^\mathrm{el}
    &=\frac{2G\nu}{1-2\nu}\cdot\frac{-2(1-2\nu)\varepsilon_0}{3(1-\nu)}
    =-\frac{4G\nu\varepsilon_0}{3(1-\nu)},\notag\\
  2Ge_{11}^\mathrm{el}
    &=-\frac{4G(1-2\nu)\varepsilon_0}{9(1-\nu)}.
    \notag
\end{align}
Adding the hydrostatic and deviatoric parts, with common denominator
$9(1-\nu)$:
\begin{align}
  \sigma_{11}^\mathrm{in}&=-\frac{12G\nu\varepsilon_0}{9(1-\nu)}
    -\frac{4G(1-2\nu)\varepsilon_0}{9(1-\nu)}\notag\\
  &=-\frac{G\varepsilon_0}{9(1-\nu)}[12\nu+4-8\nu]
    =-\frac{4G(1+\nu)\varepsilon_0}{9(1-\nu)}.
    \label{eq:sig11in}
\end{align}

\paragraph{Step~4 (energy contraction)}
By isotropy all three diagonal stress components are equal,
$\sigma_{11}^\mathrm{in}=\sigma_{22}^\mathrm{in}=\sigma_{33}^\mathrm{in}$,
so the double contraction with the diagonal eigenstrain
$\varepsilon^*_{ii}=\varepsilon_0/3$ collapses to three identical
terms: $\sigma_{ij}^\mathrm{in}\varepsilon^*_{ij}
=3\sigma_{11}^\mathrm{in}(\varepsilon_0/3)=\sigma_{11}^\mathrm{in}\varepsilon_0$.
Substituting into the energy theorem,
\begin{equation}
  \boxed{\dF^\text{dilat}=\half\beta_2 G\varepsilon_0^2 V_0,\quad
  \beta_2=\frac{4(1+\nu)}{9(1-\nu)}.}
  \label{eq:Edilat}
\end{equation}
Like $\beta_1$, the dilatational constraint factor $\beta_2$ is
strictly between 0 and 1, reflecting that the surrounding matrix
relaxes part of the volumetric eigenstrain rather than resisting it
completely; $\beta_2=0.889$ at $\nu=1/3$ is closer to unity than
$\beta_1=0.446$ because a spherical (volumetric) misfit is
accommodated less efficiently by the elastic matrix than a shear
misfit is.

\subsubsection{Total activation energy}

Shear and dilatational eigenstrains involve orthogonal tensor
components ($\varepsilon^*_{13}$ vs.\ $\varepsilon^*_{ii}$);
their cross-terms vanish upon contraction with the isotropic stiffness
tensor, so the two energies computed separately above can simply be
added:
\begin{equation}
  \boxed{\dF=\half(\beta_1\gz^2+\beta_2\varepsilon_0^2)GV_0,\quad
  V_0=\tfrac{2\pi}{3}W^3.}
  \label{eq:dF0}
\end{equation}
The two constraint factors $\beta_1$ and $\beta_2$ together encode
the resistance of the surrounding elastic matrix to the \OSTZ\ shear
and volume change respectively; both are less than~1 because the
matrix partially accommodates each eigenstrain rather than resisting
it as a rigid, undeformable cage, which is why the stored energy is
always less than that of a hypothetical free (unconstrained)
transformation.

\paragraph{Numerical check} This first-principles elastic result can
be checked directly against the empirically calibrated activation
energy used in the original Padmanabhan et al.\ grain-boundary
sliding theory. Using the canonical grain-boundary values
$\gz=0.10$, $\varepsilon_0=0.05$, $\beta_1=0.446$, $\beta_2=0.889$,
$W=2.5b=0.75$~nm (so $V_0=0.884$~nm$^3$), and $G=2.2\times10^{10}$~Pa
for Al:
$\dF=\half(0.00446+0.00222)\times2.2\times10^{10}\times0.884\times10^{-27}
=6.5\times10^{-20}$~J$=0.406$~eV, within 6\% of the calibrated value
$\approx0.38$~eV, confirming that the Eshelby-derived constraint
factors reproduce the empirically known activation energy without
any further adjustment.

\begin{table}[htbp]
\centering
\small
\caption{Canonical \OSTZ\ parameters.}
\label{tab:canonical}
\begin{tabular}{lccc}
\toprule
Parameter & Symbol & GB/superpl. & Crystal\\
\midrule
OSTZ radius & $W$ & $2.5b$ & $\approx b$\\
Shear eigenstrain & $\gz$ & 0.10 & $s_m$ (0.10--0.20)\\
Dilat.\ eigenstrain & $\varepsilon_0$ & 0.05 & ---\\
Shear constraint & $\beta_1$ & 0.446 & $\approx1$\\
Dilat.\ constraint & $\beta_2$ & 0.889 & ---\\
Activation energy & $\dF$ & 0.38~eV & ---\\
Critical number & $N_c$ & 4 & $\approx5$--$10$\\
\bottomrule
\end{tabular}

\smallskip
{\footnotesize In the crystal interior the shear eigenstrain $\gz$ is
\emph{not} a single universal constant but the material's own
\emph{ideal shear strain} $s_m$ (the relaxed engineering shear strain at
the maximum of the first-principles $\gamma$-surface), taken from
Ogata et~al.~\cite{ogata2004}: $0.137$ (Cu), $0.200$ (Al), $0.140$ (Ni),
$0.145$ (Ag), $0.105$ (Au), $\approx0.18$ (BCC Fe, W). The single value
$\gz=0.12$ quoted in earlier work is the Cu value; using it for every
metal is the source of the large discrepancies in stacking-fault energy
and Peierls stress that the per-metal ideal shear strain removes (SI,
Sec.~S2 and Fig.~S2).}
\end{table}

%%=========================================================
\section{Far-Field Stress: The Dipole Structure}
\label{sec:dipole}

Having fixed the energetics of a single \OSTZ, we now ask what it
looks like from far away, i.e.\ how it perturbs the stress in the
surrounding grain. This is the question a materials scientist asks
when picturing how one local shear event ``talks to'' its neighbours:
does it act like a point source of stress (as a dislocation does),
or is its influence short-ranged? The answer, derived below, is that
a single \OSTZ\ is a stress \emph{dipole}, not a monopole: its field
decays two powers of $r$ faster than a dislocation's. This is the
microscopic reason a lone shear-transformation or grain-boundary
sliding event cannot by itself propagate slip across a grain, and it
sets up the chain construction of \S\ref{sec:emergence} in which many
such dipoles must cooperate to reproduce a dislocation.

\subsection{Why a dipole, not a monopole}

The force monopole vanishes:
\begin{equation}
  F_k=\oint_{\partial V_0}\sigma_{ij}^* n_j\,\mathrm dS
  =C_{ijkl}\varepsilon^*_{kl}
  \underbrace{\oint_{\partial V_0}n_j\,\mathrm dS}_{=0}=0.
  \label{eq:monopole}
\end{equation}
(The surface integral of the outward normal over any closed surface
vanishes by the divergence theorem.) The leading far-field term is
therefore the \emph{force-dipole} at order $r^{-3}$.

The full multipole expansion~\cite{mura1987} is
\begin{multline}
  \sigma_{ij}(\mathbf r)=
  \underbrace{F_k\mathcal G_{ijk}}_{\text{monopole }r^{-2}}
  +\underbrace{M_{kl}\partial_l\mathcal G_{ijk}}_{\text{dipole }r^{-3}}
  \\+\underbrace{Q_{klm}\partial_l\partial_m\mathcal G_{ijk}}_{r^{-4}}
  +\cdots
  \label{eq:multipole}
\end{multline}
Since $F_k=0$, the leading nonzero term is the dipole. The moment
tensor~\cite{mura1987} is
\begin{equation}
  M_{kl}=C_{ijkl}\varepsilon^*_{ij}V_0\;\Rightarrow\;
  M_{13}=M_{31}=G\gz V_0,
  \label{eq:moment}
\end{equation}
describing a double couple: two force pairs of equal and opposite sign,
separated by the OSTZ diameter $2W$.

\subsection{Step A: Kelvin Green's function and its derivative}
\label{sec:stepA}

The elastic displacement due to a unit point force at the origin in an
infinite isotropic medium (Kelvin--Somigliana solution~\cite{mura1987}):
\begin{equation}
  \mathcal G_{ij}(\mathbf r)=A\!\left[\frac{(3-4\nu)\delta_{ij}}{r}
  +\frac{x_ix_j}{r^3}\right],\quad A=\frac{1}{16\pi G(1-\nu)}.
  \label{eq:Gij}
\end{equation}

\paragraph{Differentiating the first term}
Since $\partial_l(1/r)=-x_l/r^3$:
\[
  \partial_l\!\left[\frac{(3-4\nu)\delta_{ij}}{r}\right]
  =-\frac{(3-4\nu)\delta_{ij}x_l}{r^3}.
\]

\paragraph{Differentiating the second term (product rule)}
$\partial x_i/\partial x_l=\delta_{il}$, $\partial(r^{-3})/\partial x_l=-3x_l/r^5$:
\[
  \partial_l\!\left(\frac{x_ix_j}{r^3}\right)
  =\frac{\delta_{il}x_j+\delta_{jl}x_i}{r^3}-\frac{3x_lx_ix_j}{r^5}.
\]

\paragraph{Combined result}
\begin{align}
  \frac{\partial\mathcal G_{ij}}{\partial x_l}
  =A\Bigl[&-\frac{(3-4\nu)\delta_{ij}x_l}{r^3}
    +\frac{\delta_{il}x_j+\delta_{jl}x_i}{r^3}
    \notag\\
  &-\frac{3x_lx_ix_j}{r^5}\Bigr].
  \label{eq:dG}
\end{align}

\subsection{Step B: Displacement from the moment tensor}
\label{sec:stepB}

For an Eshelby eigenstrain inclusion, the far-field exterior
displacement is~\cite{eshelby1957,mura1987} $u_i=-M_{kl}\partial_l\mathcal G_{ik}$.
With $M_{13}=M_{31}=G\gz V_0$ (all other components zero),
the only nonzero contributions are $(k,l)=(1,3)$ and $(3,1)$.

Substituting Eq.~\eqref{eq:dG} (with $\delta_{13}=0$):
\begin{align}
  \frac{\partial\mathcal G_{i1}}{\partial x_3}
  &=A\!\left[-\frac{(3-4\nu)\delta_{i1}x_3}{r^3}
    +\frac{\delta_{i3}x_1}{r^3}-\frac{3x_3x_ix_1}{r^5}\right],
    \notag\\
  \frac{\partial\mathcal G_{i3}}{\partial x_1}
  &=A\!\left[-\frac{(3-4\nu)\delta_{i3}x_1}{r^3}
    +\frac{\delta_{i1}x_3}{r^3}-\frac{3x_1x_ix_3}{r^5}\right].
    \notag
\end{align}

Adding term by term, the $\delta_{i1}x_3$ coefficient is
$1-(3-4\nu)=4\nu-2=2(2\nu-1)$ (same for $\delta_{i3}x_1$),
and the $r^{-5}$ terms add to give $-6Ax_1x_3x_i/r^5$:
\begin{multline}
  u_i=-G\gz V_0 A\Bigl[\frac{2(2\nu-1)(\delta_{i1}x_3+\delta_{i3}x_1)}{r^3}
  \\-\frac{6x_1x_3x_i}{r^5}\Bigr].
  \label{eq:ui}
\end{multline}

Components relevant to $\sigma_{13}$:
\begin{align}
  u_1&=-G\gz V_0 A\!\left[\frac{2(2\nu-1)x_3}{r^3}
    -\frac{6x_1^2x_3}{r^5}\right],
    \label{eq:u1}\\
  u_3&=-G\gz V_0 A\!\left[\frac{2(2\nu-1)x_1}{r^3}
    -\frac{6x_1x_3^2}{r^5}\right].
    \label{eq:u3}
\end{align}

\subsection{Step C: Stress from displacement gradients}
\label{sec:stepC}

Step B gave the displacement field $u_i(\mathbf r)$ produced by the
\OSTZ\ dipole; what we actually need for the far-field stress
(Eq.~\eqref{eq:dipole} below) is the shear stress $\sigma_{13}$, which
by Hooke's law is built from the symmetric displacement gradient
$\partial_3u_1+\partial_1u_3$. Because the shear component has
$\delta_{13}=0$, the $\lambda$-term (the dilatational part of Hooke's
law) drops out automatically, leaving a pure shear-modulus relation:
$\sigma_{13}=G(\partial_3u_1+\partial_1u_3)$. The remainder of this
subsection differentiates the two terms of $u_1$ and $u_3$ found in
Step B, term by term, and adds the results.

\paragraph{Auxiliary results}
$\partial_3(x_3/r^3)=(r^2-3x_3^2)/r^5$;
$\partial_3(x_3/r^5)=(r^2-5x_3^2)/r^7$; likewise for $x_1$.

\paragraph{Computing $\partial_3u_1$}
Differentiating Eq.~\eqref{eq:u1} with respect to $x_3$:
\begin{align}
  \partial_3u_1&=-G\gz V_0 A\Bigl[\frac{2(2\nu-1)(r^2-3x_3^2)}{r^5}
    \notag\\
  &\phantom{=-G\gz V_0 A\Bigl[}\;-\frac{6x_1^2(r^2-5x_3^2)}{r^7}\Bigr].
    \label{eq:du1dx3}
\end{align}

\paragraph{Computing $\partial_1u_3$}
By the same procedure (swap $x_1\leftrightarrow x_3$):
\begin{align}
  \partial_1u_3&=-G\gz V_0 A\Bigl[\frac{2(2\nu-1)(r^2-3x_1^2)}{r^5}
    \notag\\
  &\phantom{=-G\gz V_0 A\Bigl[}\;-\frac{6x_3^2(r^2-5x_1^2)}{r^7}\Bigr].
    \label{eq:du3dx1}
\end{align}

\paragraph{Summing and converting to unit vectors $\hat r_i=x_i/r$}
$1/r^5$ group:
$2(2\nu-1)[2-3(\hat r_1^2+\hat r_3^2)]/r^3$.
$1/r^7$ group, expanding $(x_1^2+x_3^2)r^2-10x_1^2x_3^2$:
$[-6(\hat r_1^2+\hat r_3^2)+60\hat r_1^2\hat r_3^2]/r^3$.
Collecting the $(\hat r_1^2+\hat r_3^2)$ terms:
$-6(2\nu-1)-6=-12\nu$ (verify: $12\nu-6+6=12\nu$). Hence:
\begin{multline}
  \partial_3u_1+\partial_1u_3=
  -G\gz V_0 A\cdot\frac{1}{r^3}\\
  \times\bigl[4(2\nu-1)-12\nu(\hat r_1^2+\hat r_3^2)
  +60\hat r_1^2\hat r_3^2\bigr].
  \label{eq:gradsum}
\end{multline}

\subsection{Step D: Angular tensor and far-field formula}
\label{sec:stepD}

Step C produced $\partial_3u_1+\partial_1u_3$ as a function of
position; multiplying by $G$ converts it into the physical shear
stress $\sigma_{13}$, and substituting the explicit prefactor
$A=1/[16\pi G(1-\nu)]$ reduces the result to a single radial factor
$1/r^3$ times a dimensionless function of direction only. That
dimensionless function, denoted $\mathcal T_{13}(\hat{\mathbf r})$
below, is the \emph{angular tensor}: it encodes how strongly the
\OSTZ\ promotes shear stress as a function of the direction
$\hat{\mathbf r}$ from the source, independent of distance, and its
sign and zeros (examined in Step E) determine where the \OSTZ\
assists or opposes further slip. Multiplying Eq.~\eqref{eq:gradsum}
by $G$ and inserting $A=1/[16\pi G(1-\nu)]$:
\[
  \sigma_{13}=\frac{G\gz V_0}{16\pi(1-\nu)r^3}
  \bigl[-4(2\nu-1)+12\nu(\hat r_1^2+\hat r_3^2)-60\hat r_1^2\hat r_3^2
  \bigr].
\]
Writing this as $G\gz V_0\mathcal T_{13}/[2\pi(1-\nu)r^3]$ requires
dividing the bracket by 8 (ratio of prefactors $2\pi/16\pi=1/8$):
\begin{equation}
  \boxed{\mathcal T_{13}(\hat{\mathbf r})
  =\tfrac{1}{2}-\nu+\tfrac{3\nu}{2}(\hat r_1^2+\hat r_3^2)
  -\tfrac{15}{2}\hat r_1^2\hat r_3^2.}
  \label{eq:Tij}
\end{equation}
\begin{equation}
  \boxed{\sigma_{13}(\mathbf r)=\frac{G\gz V_0}{2\pi(1-\nu)}
  \cdot\frac{\mathcal T_{13}(\hat{\mathbf r})}{r^3}.}
  \label{eq:dipole}
\end{equation}
The $r^{-3}$ decay follows from differentiating the Kelvin $r^{-1}$
Green's function twice: two orders faster than the Volterra $r^{-1}$
field.

\subsection{Step E: Verification in special directions}
\label{sec:stepE}

Equation~\eqref{eq:dipole} was derived through a long chain of
differentiations, so before using it we check it in three special
directions where the answer can be anticipated on physical grounds:
the glide plane itself (where the formula must reduce to the
expression used later in \S\ref{sec:glideplane} and
\S\ref{sec:lorentzian} to build up dislocation chains), the
boundary-normal axis (where axial symmetry of the oblate spheroid
demands a specific relation between directions), and the question of
whether $\mathcal T_{13}$ ever changes sign (since a sign change
would mean the \OSTZ\ opposes slip in some directions, undermining
its role as a cooperative nucleus).

\paragraph{Glide plane ($x_3=0$, $\hat r_3=0$)}
$\mathcal T_{13}|_{\hat r_3=0}=\half-\nu+(3\nu/2)\hat r_1^2$.
Expanding $r^2=x_1^2+x_2^2$ and collecting:
\begin{equation}
  \sigma_{13}(x_1,x_2,0)=\frac{G\gz V_0}{4\pi(1-\nu)r^5}
  [(1+\nu)x_1^2+(1-2\nu)x_2^2].
  \label{eq:glidecheck}
\end{equation}
Positive definite for $0<\nu<\half$.

\paragraph{On-axis ($\hat r_1=0$, $\hat r_3=1$)}
$\mathcal T_{13}=\half-\nu+3\nu/2=(1+\nu)/2>0$. By oblate axial
symmetry, the same result holds along $\hat x_1$.

\paragraph{Nodal surface}
Setting $\mathcal T_{13}=0$ yields no real solution for
$0<\nu<\half$: there is no nodal surface in the half-space
$x_3\ge0$. The \OSTZ\ promotes shear in all glide-plane directions
(cooperative nucleus).

\subsection{Glide-plane stress: six-step derivation}
\label{sec:glideplane}

Step E confirmed the general 3D formula in special directions, but
the calculation that matters for everything that follows, in
particular the chain construction of \S\ref{sec:emergence}, is the
stress restricted to the glide plane itself, since that is where
neighbouring \OSTZ s actually sit. Rather than simply substituting
$x_3=0$ into Eq.~\eqref{eq:dipole} and Eq.~\eqref{eq:Tij} (which is
algebraically awkward because of the way $\hat r_1$ and $\hat r_3$
enter), it is more transparent to recompute $\sigma_{13}|_{z=0}$
directly from the Green's function components, working through the
off-diagonal and diagonal contributions one at a time and only
afterwards converting to the unit-vector form. The six steps below
carry this out and then re-express the result in terms of the
effective Burgers vector $\beff$ (\S\ref{sec:beff}) so that it can be
used directly as the building block for an \OSTZ\ chain.
Working at $x_3=0$ ($r^2=x^2+y^2$, $x\equiv x_1$, $y\equiv x_2$).

\paragraph{Step~1 (off-diagonal $\mathcal G_{13}$)}
$\partial^2(\mathcal G_{13})/(\partial x_1\partial x_3)$ at $z=0$:
gives $(y^2-2x^2)/r^5$ (using $\partial_3(x_1x_3/r^3)|_{z=0}
=x_1/r^3$, then $\partial_1$).

\paragraph{Step~2 (diagonal $\mathcal G_{33}$)}
At $z=0$: $\mathcal G_{33}|_{z=0}=A(3-4\nu)/r$.
$\partial_1^2(1/r)=(3x^2-r^2)/r^5=(2x^2-y^2)/r^5$:
\[
  \partial_1^2\mathcal G_{33}\big|_{z=0}=A(3-4\nu)(2x^2-y^2)/r^5.
\]

\paragraph{Step~3 (diagonal $\mathcal G_{11}$)}
At $z=0$: $\mathcal G_{11}|_{z=0}=A[(3-4\nu)/r+x^2/r^3]$.
$\partial_3^2(1/r)|_{z=0}=-1/r^3$;
$\partial_3^2(1/r^3)|_{z=0}=-3/r^5$:
\[
  \partial_3^2\mathcal G_{11}\big|_{z=0}=A[-(3-4\nu)/r^3-3x^2/r^5].
\]
Expanding: $-(3-4\nu)r^2-3x^2=-(6-4\nu)x^2-(3-4\nu)y^2$.

\paragraph{Step~4 (collect numerator, units of $A/r^5$)}
\begin{align*}
  x^2\text{ coeff.:}&\;2(y^2\!-\!2x^2)\to-4\\
  &+(6\!-\!8\nu)-(3\!-\!4\nu)-3=-4(1+\nu),\\
  y^2\text{ coeff.:}&\;2-(3\!-\!4\nu)-(3\!-\!4\nu)=-4(1-2\nu).
\end{align*}
Numerator: $-4A[(1+\nu)x^2+(1-2\nu)y^2]/r^5$.

\paragraph{Step~5 (assemble)}
\[
  \sigma_{13}(x,y,0)=\frac{G\gz V_0[(1+\nu)x^2+(1-2\nu)y^2]}
  {4\pi(1-\nu)r^5}.
\]

\paragraph{Step~6 (re-express with $\beff$, restore finite $W$)}
Using $\beff=\gz W$:
\begin{align*}
  G\gz V_0&=G\gz\,\tfrac{2\pi W^3}{3}=G\beff\,\tfrac{2\pi W^2}{3},\\
  \text{so}\quad\frac{G\gz V_0}{4\pi(1-\nu)}&=\frac{G\beff W^2}{6(1-\nu)}.
\end{align*}
Regularising $r^2\to r^2+W^2$:
\begin{equation}
  \boxed{\sigma_{13}(x,y,0)=
  \frac{G\beff W^2[(1+\nu)x^2+(1-2\nu)y^2]}
  {6(1-\nu)(x^2+y^2+W^2)^{5/2}}.}
  \label{eq:glideplane}
\end{equation}
Positive definite; finite at the origin.

\paragraph{Physical interpretation}
Along $\hat x$ ($y=0$): promotes forward shear.
Along $\hat y$ ($x=0$): promotes lateral shear.
No nodal planes in the glide plane: the \OSTZ\ is a
\emph{cooperative nucleus}.

\subsection{Effective Burgers vector}
\label{sec:beff}

The glide-plane formula just derived, Eq.~\eqref{eq:glideplane},
already anticipated that the natural strength parameter for an
\OSTZ\ is not $\gz$ alone but a length-scaled combination $\beff$
with the same units as a dislocation Burgers vector; this is the
quantity that will later sum, $N$ at a time, into the total Burgers
vector of an \OSTZ\ chain (\S\ref{sec:PN}). Three independent routes
to $\beff$ are checked against each other below; they agree to
within $O(1)$, and the discrepancy itself is informative about the
approximations involved.

\paragraph{Method~1 (kinematic, exact)} The eigenstrain
$\varepsilon^*_{13}=\gz/2$ is, by definition, half the engineering
shear strain $\gz=\partial u_1/\partial x_3$ relating the relative
glide displacement $u_1$ to the through-thickness coordinate $x_3$.
The total relative displacement between the two flat faces of the
spheroid, separated by the full thickness $2a_3=W$, is obtained by
integrating this constant strain across the thickness:
\[
  \beff=\int_{-a_3}^{a_3}\gz\,\mathrm dx_3=\gz\times2a_3=\gz W.
\]
This is an exact kinematic statement, independent of the elastic
solution, and is adopted as the primary definition.

\paragraph{Method~2 (moment tensor)} As an independent check, $\beff$
can instead be extracted from the seismic moment $M_{13}=G\gz V_0$
of Eq.~\eqref{eq:moment} by treating the \OSTZ\ as a thin disk source
of area $\pi W^2$ carrying a uniform Burgers vector $\beff$, so that
$M_{13}=G\beff\,\pi W^2$. Equating the two expressions for $M_{13}$
and substituting $V_0=(2\pi/3)W^3$:
\begin{align*}
  G\beff\pi W^2&=G\gz V_0=G\gz\cdot\frac{2\pi W^3}{3}\\
  \Rightarrow\quad\beff&=\frac{\gz\cdot\frac{2\pi W^3}{3}}{\pi W^2}=\frac{2\gz W}{3},
\end{align*}
which underestimates Method~1 by 33\%.
\paragraph{Origin of the discrepancy} The seismic-moment formula
equates $M_{13}$ to a Burgers vector spread uniformly over the disk
area $\pi W^2$, using the \emph{full} oblate-spheroid volume
$V_0=(2\pi/3)W^3$ rather than the actual projected slip area. For a
thin oblate disk ($\alpha\ll1$) Methods~1 and~2 converge; at
$\alpha=1/2$ the 33\% gap is non-negligible and is carried forward
as an $O(1)$ uncertainty rather than resolved.

\paragraph{Method~3 (Lorentzian normalisation)} The third, weaker check
comes from the Lorentzian dislocation density derived later in
\S\ref{sec:lorentzian}, $\rho_1(x)=(\beff/\pi)W/(x^2+W^2)$
(Eq.~\eqref{eq:rho1}), which integrates to exactly $\beff$ over all
$x$. This, however, holds \emph{by construction}: the Lorentzian's
prefactor is chosen precisely so that the normalisation comes out to
$\beff$, so Method~3 cannot serve as an independent confirmation of
the value found in Methods~1 and~2; it merely shows that the
later construction is self-consistent with whichever $\beff$ is fed
into it.

Method~1 is adopted as the primary definition throughout the rest of
the paper:
\begin{equation}
  \boxed{\beff=\gz W.}
  \label{eq:beff}
\end{equation}

\begin{table}[htbp]
\centering
\small
\caption{Single \OSTZ\ vs.\ Volterra dislocation.}
\label{tab:compare1}
\renewcommand{\arraystretch}{1.35}
\begin{tabular}{p{1.6cm}p{2.3cm}p{2.3cm}}
\toprule
Property & Single \OSTZ & Volterra\\
\midrule
Stress decay & $r^{-3}$ (dipole) & $r^{-1}$ (monopole)\\
Topo.\ charge & 0 (neutral) & $\pm b$\\
Net $b$ & 0 & $b$\\
Core stress & Finite ($\le G\gz$) & $\to\infty$\\
Influence & $\sim W$ & $\infty$\\
Elastic $E$ & $\half\beta_1\gz^2GV_0$ & $Gb^2L\ln(R/r_0)$\\
Thermal act. & Intrinsic & Appended\\
\bottomrule
\end{tabular}
\renewcommand{\arraystretch}{1}
\end{table}

\subsection{Why a single \OSTZ\ cannot drive macroscopic slip}
\label{sec:nomacroslip}

The dipole structure ($r^{-3}$ decay, Eq.~\eqref{eq:dipole}) has a
decisive consequence: a single \OSTZ\ cannot drive long-range slip.
Its stress influence radius is effectively $\sim W$; beyond $\sim5W$
the stress has fallen below 1\% of its peak value. Only when
$N=N_c$ \OSTZ s cooperate does the collective Burgers circuit close
on a quantised $b\in\Lambda_\mathrm{Bravais}$
(\S\ref{sec:Nc}), enabling the long-range $r^{-1}$ field of a full
dislocation (Eq.~\eqref{eq:volterra}). Table~\ref{tab:compare1}
summarises the contrast.

%%=========================================================
\section{Emergence of Dislocations from \OSTZ\ Chains}
\label{sec:emergence}

This section answers the question raised in
\S\ref{sec:nomacroslip}: how does a row of short-ranged dipoles add
up to the long-ranged, singular field of a textbook dislocation? We
proceed in the same way a materials scientist would build up a
slipped region from individual atomic-scale events: line up $N$
\OSTZ s along the glide direction, sum their fields, and ask what
collective object emerges. The result is exact and structurally
important: an aligned \OSTZ\ chain is mathematically identical to a
\PN\ dislocation, with the \OSTZ\ radius $W$ playing the role of the
phenomenological core width that is normally fitted to experiment.

\subsection{Lorentzian dislocation density: four steps}
\label{sec:lorentzian}

\subsubsection{Step A: 2D dislocation density (BCS compliance)}
\label{sec:stepA-BCS}

The Bilby--Cottrell--Swinden (BCS) compliance~\cite{bilby1963}
$\delta u_1=2(1-\nu)\sigma_{13}/G$ relates local slip to traction.
\paragraph{Physical basis} A shear traction $\sigma_{13}$ on the
glide plane produces a local slip increment $\delta u_1$ smaller
than the naive estimate $\sigma_{13}/G$ by a factor $(1-\nu)$,
because the surrounding half-space elastically constrains the
opening of the slip increment in Poisson-ratio fashion; the factor
can be obtained from the J-integral analysis of a shear
crack~\cite{rice1968} or, equivalently, from the dislocation-density
relation $\sigma=G\rho/[2(1-\nu)]$ of the \PN\
model~\cite{hirth1982}.
\paragraph{Applicability to the \OSTZ\ $r^{-3}$ field} Three features
justify using the planar BCS compliance here even though it was
derived for smooth slip distributions: (i)~the \OSTZ\ field decays
as $r^{-3}$, faster than the dislocation's $r^{-1}$, so the
compliance integral converges absolutely without a short-distance
cutoff; (ii)~for $N\gg1$ \OSTZ s the collective field is smooth on
scales $\gg W$ and approaches the \PN\ continuum limit; (iii)~for
grain-boundary \OSTZ s ($W=2.5b$) the continuum approximation is
well justified, while for crystal-interior \OSTZ s ($W\approx b$) it
is marginal and atomistic corrections of $O(b/W)$ may be
non-negligible.
\paragraph{Why an isotropic approximation is needed} The exact
glide-plane stress derived in \S\ref{sec:glideplane},
Eq.~\eqref{eq:glideplane}, is not a convenient starting point for
the Abel projection of Step B below, for two reasons. First, along
the chain axis itself ($y=0$) it reduces to
$\sigma_{13}(x,0,0)\propto(1+\nu)x^2/(x^2+W^2)^{5/2}$, which
\emph{vanishes at $x=0$} and only rises to a maximum away from the
origin, an awkward, non-monotonic shape to seed a dislocation-density
profile that should peak at the centre of the \OSTZ. Second, its
angular dependence on $(1+\nu)x^2+(1-2\nu)y^2$ makes the $y$-integral
of Step B intractable in closed form. We therefore replace the exact
field by the simplest possible substitute that preserves its two
physically essential features, the correct $r^{-3}$ far-field decay
and positivity everywhere in the glide plane, while being maximal
(rather than zero) at the origin: a spherically symmetric profile
with the same prefactor and the same regularisation length $W$,
\begin{equation}
  \sigma_{13}(x,y,0)\approx
  \frac{G\beff W^2}{2\pi(1-\nu)(x^2+y^2+W^2)^{3/2}}.
  \label{eq:gpisotropic}
\end{equation}
This is an explicit modelling choice, not a re-derivation of
Eq.~\eqref{eq:glideplane}: its justification is that it reproduces
the known asymptotics and, as shown below, integrates to exactly the
Lorentzian dislocation density that both \OSET\ and the independent
\PN\ variational calculation select (\S\ref{sec:stepC-lorentz}).

With Eq.~\eqref{eq:gpisotropic} in hand, the 2D dislocation density
follows by substituting it into the BCS compliance relation
$\rho_{2D}=\delta u_1=2(1-\nu)\sigma_{13}/G$ derived above:
\begin{equation}
  \rho_{2D}(x,y)=\frac{2(1-\nu)}{G}\cdot
  \frac{G\beff W^2}{2\pi(1-\nu)(x^2+y^2+W^2)^{3/2}}.
  \label{eq:rho2Dsub}
\end{equation}
The factors of $G$ and $2(1-\nu)$ cancel between the compliance
prefactor and $\sigma_{13}$, leaving
\begin{equation}
  \rho_{2D}(x,y)=\frac{\beff W^2}{\pi(x^2+y^2+W^2)^{3/2}}.
  \label{eq:rho2D}
\end{equation}

\paragraph{Normalisation check}
Converting to polar coordinates $r^2=x^2+y^2$:
\begin{align*}
  \iint\rho_{2D}\,\mathrm dx\,\mathrm dy
  &=\frac{\beff W^2}{\pi}\cdot2\pi
    \int_0^\infty\frac{r\,\mathrm dr}{(r^2+W^2)^{3/2}}.
\end{align*}
Substitution $u=r^2+W^2$, $du=2r\,dr$:
$\int=\half[-2/\sqrt{u}]_{W^2}^\infty=1/W$.
Therefore $\iint\rho_{2D}=2\beff W$.
(The result is $2\beff W$, not $\beff$, because $\rho_{2D}$ is slip
per unit glide-plane area; the Abel projection recovers the correct
normalisation below.)

\subsubsection{Step B: Abel projection}
\label{sec:stepB-abel}

Integrate over $y$ with $a^2=x^2+W^2$; let $y=a\tan\theta$:
\begin{equation}
  \int_{-\infty}^\infty\frac{\mathrm dy}{(a^2+y^2)^{3/2}}
  =\frac{1}{a^2}\!\int_{-\pi/2}^{\pi/2}\cos\theta\,\mathrm d\theta
  =\frac{2}{x^2+W^2}.
  \label{eq:abel}
\end{equation}
Raw 1D density: $\rho_1^\text{raw}=2\beff W^2/[\pi(x^2+W^2)]$,
which integrates to $2\beff W$.

\subsubsection{Step C: Lorentzian (normalised)}
\label{sec:stepC-lorentz}

Divide by $2W$ to enforce $\int\rho_1\,\mathrm dx=\beff$:
\begin{equation}
  \boxed{\rho_1(x)=\frac{\beff}{\pi}\cdot\frac{W}{x^2+W^2}.}
  \label{eq:rho1}
\end{equation}
\paragraph{Check}
$\int_{-\infty}^\infty\rho_1\,\mathrm dx=(\beff W/\pi)(\pi/W)=\beff$.
$\square$

The Lorentzian is the unique non-negative distribution integrating to
$\beff$ with width $W$ and reducing to a Dirac delta as $W\to0$.
Both \OSET\ and the P--N model are governed by the same
Cauchy--Hilbert Green's operator of planar
elasticity~\cite{muskhelishvili1953,bilby1963}: the singular
integral kernel $1/(x-x')$ that relates a planar dislocation
density to the shear stress it produces (Eq.~\eqref{eq:conv} below)
is the Cauchy kernel of the Hilbert transform, and the Lorentzian is
precisely the density for which that transform reproduces the
non-singular, finite-core stress field of \S\ref{sec:chainStress}.

\subsubsection{Step D: Fourier transform}
\label{sec:stepD-fourier}

For $k>0$, close the contour in the lower half-plane (where
$e^{-ikx}\to0$). The only enclosed pole is at $x=-iW$:
\[
  \mathrm{Res}\!\left[\frac{We^{-ikx}}{\pi(x^2+W^2)},\,-iW\right]
  =\frac{e^{-kW}}{-2iW}.
\]
Clockwise contour ($-2\pi i\times\text{Res}$):
\begin{equation}
  \hat\rho(k)=\beff\,e^{-|k|W}.
  \label{eq:rhohat}
\end{equation}
At $k=2\pi/b$: $\hat\rho=\beff e^{-2\pi W/b}$ (Peierls lattice-filter
factor). The OSTZ core acts as a low-pass filter: high-spatial-frequency
lattice components are exponentially suppressed.

\subsection{Theorem: the $N$-\OSTZ\ chain is a \PN\ dislocation}
\label{sec:PN}

We can now state precisely what a chain of \OSTZ s is, in dislocation
language. Consider $N$ \OSTZ s packed densely along the glide
direction, each contributing its own Lorentzian density
$\rho_1(x)$, Eq.~\eqref{eq:rho1}, centred on its own position; in the
dense-packing limit the individual Lorentzians simply superpose
(their characteristic width $W$ being the spacing scale), giving a
single chain density $N$ times larger than that of one \OSTZ:
$\rho_N(x)=N\cdot(\beff/\pi)W/(x^2+W^2)$.

\paragraph{Theorem (dislocation emergence)} A co-planar chain of $N$
\OSTZ s with aligned shear eigenstrains is mathematically identical
to a \PN\ dislocation with total Burgers vector $b=N\beff$, core
half-width $\zeta=W$, and core energy $E_\text{core}=N\dF$.

\paragraph{Proof} The cumulative slip $U(x)$ at a field point $x$ is,
by definition, the total Burgers vector already carried by every
\OSTZ\ to the left of $x$, i.e.\ the integral of the chain density
from $-\infty$ up to $x$:
\begin{align}
  U(x)&=\int_{-\infty}^x\rho_N(x')\,\mathrm dx'\notag\\
  &=\frac{N\beff W}{\pi}\cdot\frac{1}{W}
    \!\left[\arctan\frac{x'}{W}\right]_{-\infty}^x
    \notag\\
  &=\frac{N\beff}{2}+\frac{N\beff}{\pi}\arctan\frac{x}{W}.
  \label{eq:slip}
\end{align}
(Using $\arctan(-\infty)=-\pi/2$.)

Setting $b\equiv N\beff$ and $\zeta\equiv W$:
\begin{equation}
  U(x)=\frac{b}{2}+\frac{b}{\pi}\arctan\frac{x}{\zeta}.
  \label{eq:PNprofile}
\end{equation}
This is exactly the \PN\ displacement profile~\cite{peierls1940,nabarro1947},
obtained there from an entirely different starting point, the
variational minimisation of elastic self-energy against a sinusoidal
misfit potential. \textit{Limits:} as $x\to-\infty$, $U\to0$ (no slip
has yet occurred); as $x\to+\infty$, $U\to b$ (the full chain
Burgers vector has passed); at $x=0$, $U=b/2$, the midpoint of the
transition. The continuum dislocation density is the spatial
derivative $\rho(x)=\partial U/\partial x$; using
$\mathrm d/\mathrm dx[\arctan(x/W)]=W/(x^2+W^2)$,
\[
  \rho(x)=\frac{N\beff}{\pi}\cdot\frac{W}{x^2+W^2}=\rho_N(x),
\]
which reproduces the chain density we started from, confirming
self-consistency. $\square$

\paragraph{Corollary} The core half-width $\zeta=W$ is derived from
\OSTZ\ geometry, not a free parameter, and the core energy
$E_\text{core}=N\dF$ is simply the accumulated activation energy of
the $N$ \OSTZ s that make up the chain (quantified explicitly in
\S\ref{sec:coreE}). The \PN\ model is therefore not an independent
theory: it is the continuum limit of the \OSTZ\ ensemble. It is
worth noting that the commonly used algebraic approximation to the
\PN\ profile, $u_\text{approx}(x)=\tfrac{b}{2}[1+x/\sqrt{x^2+\zeta^2}]$,
is \emph{not} reproduced here: its slope at the origin,
$b/2\zeta$, differs from the exact slope $b/\pi\zeta$ of
Eq.~\eqref{eq:PNprofile} by a factor $\pi/2\approx1.57$, so the
\OSTZ\ chain construction selects the exact arctangent profile
rather than this common approximation.

\subsection{Chain stress field and Volterra limit}
\label{sec:chainStress}

Having shown that the \emph{displacement} field of the \OSTZ\ chain
is exactly the \PN\ profile, we now compute its \emph{stress} field
directly, both to verify consistency with the \PN\ result and
because the stress, not the displacement, is what determines whether
two dislocation segments attract, repel, or recombine in later
applications. A planar distribution of slip $\rho_N(x')$ generates a
shear stress through the same Cauchy--Hilbert kernel introduced in
\S\ref{sec:lorentzian}: each element of slip $\rho_N(x')\mathrm dx'$
at $x'$ contributes a stress $\propto1/(x-x')$ at the field point
$x$, and summing (integrating) these contributions over the whole
chain gives
\begin{equation}
  \sigma_{13}^{(N)}(x)=\frac{G}{2\pi(1-\nu)}\,\mathrm{P.V.}
  \int_{-\infty}^\infty\frac{\rho_N(x')}{x-x'}\,\mathrm dx'.
  \label{eq:conv}
\end{equation}
The integral must be read as a Cauchy principal value because the
kernel $1/(x-x')$ has a pole at $x'=x$; physically, this reflects the
fact that the stress immediately at a point of the glide plane is
not affected by the (symmetric, cancelling) contributions of slip
infinitesimally to either side of it.

\paragraph{Partial fractions}
The integral is evaluated by separating the awkward product of
denominators $(x'^2+W^2)(x-x')$ into two simpler pieces, each of
which can be integrated by elementary means. Decompose
$W/[(x'^2+W^2)(x-x')]$:
$W=A(x'^2+W^2)+(Bx'+C)(x-x')$.
Matching powers of $x'$:
$x'^2$: $A=B$;
$x'^1$: $Bx-C=0\Rightarrow C=Ax$;
constant: $A(x^2+W^2)=W\Rightarrow A=W/(x^2+W^2)$.
Hence:
\[
  \frac{W}{(x'^2+W^2)(x-x')}
  =\frac{W}{x^2+W^2}\!\left[\frac{1}{x-x'}
  +\frac{x'+x}{x'^2+W^2}\right].
\]
This split isolates the singular Cauchy-kernel piece (Term~1) from a
smooth, elementary integral (Term~2).

\paragraph{Term~1}
P.V.\,$\int\mathrm dx'/(x-x')=0$ (integrand is odd about $x'=x$): the
principal-value prescription removes exactly this divergent,
antisymmetric contribution, leaving no trace of the pole in the
final answer.

\paragraph{Term~2}
$\int x'/(x'^2+W^2)\,\mathrm dx'=0$ (odd in $x'$).
$\int\mathrm dx'/(x'^2+W^2)=\pi/W$
(via $x'=W\tan\theta$).
So Term~2 gives $\pi x/(x^2+W^2)$.

\paragraph{Combining}
Substituting both terms back into Eq.~\eqref{eq:conv} and
multiplying through by the prefactor $G/[2\pi(1-\nu)]$ and the
chain's total Burgers vector $N\beff$ (carried by $\rho_N$),
\begin{equation}
  \boxed{\sigma_{13}^{(N)}(x)=\frac{GN\beff}{2\pi(1-\nu)}
  \cdot\frac{x}{x^2+W^2}.}
  \label{eq:chainStress}
\end{equation}
This stress is non-singular at $x=0$ (it vanishes there, by the
antisymmetry already used in Term~1) and reaches its maximum at
$x=W$, i.e.\ at the edge of the core rather than at its centre,
mirroring the familiar shape of the classical \PN\ stress profile
but with $W$, rather than a fitted core radius, fixing its width.

\paragraph{Volterra limit ($W\to0$)}
The chain stress of Eq.~\eqref{eq:chainStress} was derived for finite
core width $W$ and is everywhere finite; the classical singular
Volterra field is recovered only as a limiting case, obtained by
shrinking the core to zero while holding the total Burgers vector
$b=N\beff$ fixed (so $N\to\infty$ as $W\to0$):
$\lim_{W\to0}x/(x^2+W^2)=1/x$, so:
\begin{equation}
  \sigma_{13}^\text{Volterra}=\frac{Gb}{2\pi(1-\nu)x}.
  \label{eq:volterra}
\end{equation}
which is precisely the classical Volterra dislocation stress field.
Three observations follow. First, the singularity is a
\emph{collective} effect: no single finite-$W$ \OSTZ\ in the chain is
itself singular, and the $1/x$ field only emerges in the joint limit
$N\to\infty$, $W\to0$ taken together. Second, this limit requires
$W\to0$, which is unphysical: the actual \OSTZ\ radius is always of
order the lattice spacing, $W\sim b$, never vanishingly small. Third,
and as a consequence of the first two points, \OSET\ contains no
physical singularity anywhere; the Volterra singularity exists only
in this idealised, never-actually-reached mathematical limit, and
the physically correct description of a real dislocation core is
always the finite-$W$ result of Eq.~\eqref{eq:chainStress}.

\paragraph{Burgers circuit}
The chain construction also makes the Burgers-vector quantisation of
classical dislocation theory transparent rather than axiomatic. For
any circuit $C$ that encloses all $N$ \OSTZ s of the chain, the
closure failure is just the total slip carried by the chain, i.e.\
the integral of $\rho_N(x)$ across the full length $L$ of the
circuit, evaluated in the $L\to\infty$ limit:
\[
  \oint_C\mathrm du_1=\int_{-L/2}^{L/2}\rho_N(x)\,\mathrm dx
  \xrightarrow{L\to\infty}N\beff=b.
\]
The Burgers vector $b=N\gz W$ is therefore simply read off the
chain's total slip; it becomes \emph{quantised} onto a Bravais
lattice vector, $b\in\Lambda_\mathrm{Bravais}$, only once the
surrounding crystal lattice forces the chain length to lock onto the
specific value $N=N_c$ derived next. In an amorphous solid, with no
lattice to enforce this locking, no such quantisation occurs and $b$
remains a continuous, unquantised quantity.

\subsection{Critical cluster number}
\label{sec:Nc}

The lattice locks the chain length once the accumulated chain
Burgers vector $N\beff$ reaches the lattice's own Burgers vector
$b_\mathrm{latt}$ exactly; setting $N\beff=b_\mathrm{latt}$ and
solving for $N$ gives the critical cluster number
\begin{equation}
  \boxed{N_c=\frac{b_\mathrm{latt}}{\gz W}.}
  \label{eq:Nc}
\end{equation}
For FCC copper ($W=b$, $\gz=0.137$) this gives $N_c\approx7$ (and
$N_c\approx5$--$10$ across the FCC and BCC metals); for the
grain-boundary regime ($W=2.5b$, $\gz=0.10$) it gives $N_c=4$ exactly,
a value confirmed experimentally by the cooperative sliding events
observed in Zn--22Al~\cite{astanin-part2}.

This single parameter also explains, without further assumptions,
why dislocations do not form in amorphous materials. In a
crystalline material the periodic lattice potential forces \OSTZ s
to cluster all the way to $N=N_c$ before a stable configuration is
reached, at which point a genuine lattice dislocation nucleates and
subsequently propagates as a rigid, quantised whole. In an amorphous
solid there is no periodic potential to enforce this clustering, so
\OSTZ s instead activate individually or in small, sub-critical
clusters with $N\ll N_c$, producing STZ-like deformation rather than
dislocation glide. A grain boundary occupies an intermediate
position: the quasi-crystalline order along the boundary plane
permits partial, but not complete, quantisation, which is why
boundary dislocations and partial dislocations, rather than full
lattice dislocations, are the natural carriers of grain-boundary
plasticity.

%%=========================================================
\section{Dislocation Quantities as Derived Results}
\label{sec:derived}

The chain construction of \S\ref{sec:emergence} would be a
mathematical curiosity unless it reproduces the handful of numbers
that every materials scientist uses to characterise dislocation
behaviour: the stress needed to shear a perfect lattice, the
Peierls stress that pins a dislocation between Peierls valleys, the
energy locked in its core, the Frank--Read stress that activates a
dislocation source, and the stacking-fault energy that controls
deformation-twinning and partial-dislocation separation. In
conventional theory each of these requires its own phenomenological
input ($r_0$, $\alpha_\text{core}$, an empirical $\gamma_\text{SF}$,
\ldots). Below, all five are obtained from the same three \OSTZ\
parameters ($W$, $\gz$, $G$) already fixed in \S\ref{sec:ostz},
with no additional fitting.

\subsection{Theoretical shear strength}
\label{sec:strength}

\paragraph{Physical argument} The characteristic activation stress
$\tau^*$ is the applied shear stress at which the mechanical work
done on an \OSTZ\ exactly equals its activation-energy barrier: the
work done by $\tau^*$ shearing the \OSTZ\ material through strain
$\gz$ over volume $V_0$ is $W_\mathrm{mech}=\tau^*\gz V_0$. Setting
$W_\mathrm{mech}=\dF$ gives the cascade instability threshold
$\tau^*\gz V_0=\dF=\half\beta_1\gz^2GV_0$:
\begin{equation}
  \boxed{\tau^*=\frac{\beta_1\gz G}{2}.}
  \label{eq:strength}
\end{equation}
For $\beta_1\approx1$ and the Cu ideal shear strain $\gz=0.137$:
$\tau^*\approx G/15$ (Frenkel range~\cite{frenkel1926} $G/10$--$G/30$).

\subsection{Peierls stress}
\label{sec:peierls}

\subsubsection{OSET-derived amplitude}

The Peierls stress is set by the amplitude of the periodic
lattice-misfit potential the dislocation core must overcome as it
glides; classically this amplitude is borrowed from the Frenkel
shear strength. Here it is instead fixed self-consistently by the
\OSTZ\ theoretical shear strength of \S\ref{sec:strength}, removing
that external input. Setting $\tau^*=(2\pi/b)V_0^\text{OSET}$:
\begin{equation}
  V_0^\text{OSET}=\frac{\beta_1\gz Gb}{4\pi}.
  \label{eq:V0OSET}
\end{equation}
For Cu: $\tau^*\approx G/15$ vs.\ the Frenkel estimate~\cite{frenkel1926} $G/[2(1-\nu)]\approx0.75G$.

\subsubsection{Fourier convolution}

With the misfit amplitude fixed, the Peierls barrier is obtained by
convolving the periodic lattice potential with the \OSTZ\ dislocation
density derived in \S\ref{sec:lorentzian}, exploiting the fact that
the Lorentzian's Fourier transform (Eq.~\eqref{eq:rhohat}) makes this
convolution elementary.
Misfit energy $E_\text{mis}=\int V(x-x_0)\rho(x)\,\mathrm dx$,
$V=V_0[1-\cos(2\pi x/b)]$.
Using $\hat\rho(k_1)=be^{-2\pi W/b}$ at $k_1=2\pi/b$
and Euler's formula $\cos=\mathrm{Re}[e^{i\theta}]$:
\begin{equation}
  E_\text{mis}(x_0)\big|_\text{osc}=
  -\frac{Gb^2}{4\pi(1-\nu)}\,e^{-2\pi W/b}
  \cos\!\left(\frac{2\pi x_0}{b}\right).
  \label{eq:Emisfit}
\end{equation}
(Constant $V_0 b$ term dropped as it plays no role in $\tau_P$.)

\subsubsection{Peierls barrier and half-space correction}

The misfit energy found above oscillates as the dislocation glides
from one lattice site to the next; the Peierls stress is the
maximum slope of that oscillation, after a geometric correction for
the fact that the glide plane is a free (or bounded) surface rather
than embedded in an infinite medium.
Barrier: $\Delta E_P=2Ae^{-2\pi W/b}$ with $A=Gb^2/[4\pi(1-\nu)]$.
Maximum slope: $\tau_P=(2\pi A/b^2)e^{-2\pi W/b}=G e^{-2\pi W/b}/[2(1-\nu)]$.
With half-space correction~\cite{peierls1940} ($W\to W/(1-\nu)$, prefactor $\times2$):
\begin{equation}
  \boxed{\tau_P=\frac{2G}{1-\nu}
  \exp\!\left(-\frac{2\pi W}{b(1-\nu)}\right).}
  \label{eq:peierls}
\end{equation}
Cu ($W=b$, $\nu=1/3$): $\tau_P\approx2.1\times10^{-4}G$
(experiment~\cite{nabarro1997}: $\sim10^{-4}G$). The width that enters
the exponential is the Peierls glide-misfit width $W_P$, the one
material-specific quantity any Peierls--Nabarro model must supply; with
the per-metal $W_P$ ($\sim1.1$--$1.6\,b$ for the wide planar FCC cores,
$\sim0.5$--$0.65\,b$ for the compact non-planar BCC screw cores) the same
formula reproduces the experimental Peierls stress for all seven FCC and
BCC metals (SI, Fig.~S2). The earlier choice $W_P=b$ for every metal made
the Al value $\sim200\times$ too high and the BCC values $\sim100\times$
too low.

\subsection{Core energy}
\label{sec:coreE}

In classical theory the dislocation core energy is a fitted
constant, $E_\text{core}=\alpha Gb^2$ with $\alpha\approx0.1$--$1$
chosen to match experiment or atomistic simulation, because the
elastic continuum solution diverges at the core and contributes
nothing there. In \OSET\ the core is simply the $N_c$ \OSTZ s that
have not yet collapsed into the far-field Volterra line, so its
energy is just their accumulated activation energy:
$E_\text{core}=N_c\dF$. Substituting $N_c=b/(\gz W)$ and
$\dF=\pi\beta_1\gz^2GW^3/3$:
\begin{equation}
  \boxed{E_\text{core}=\frac{b}{\gz W}\cdot\frac{\pi\beta_1\gz^2GW^3}{3}
  =\frac{\pi\beta_1 b\gz GW^2}{3}.}
  \label{eq:coreE}
\end{equation}
Writing $E_\text{core}=\alpha Gb^2$: for Cu ($W\approx b$, $\gz=0.137$),
$\alpha=\pi\beta_1\gz/3\approx0.14$ (exp/DFT~\cite{woodward2002}: $0.1$--$0.2$).
With the per-metal ideal shear strain and the structural width $W=b$
(distinct from the narrow Peierls width $W_P$ above), the FCC core
energies fall in the DFT range and the BCC values within $\sim1.6\times$
(SI, Fig.~S2).

\subsection{Frank--Read critical stress}
\label{sec:FR}

Following Frank \& Read~\cite{frank-read1950}:

\paragraph{Step~1 (line tension)}
The line tension $T$ is the dislocation's elastic self-energy per
unit length, obtained by integrating the energy density of the
$\sigma\sim Gb/r$ stress field over the annular region
$r_0<r<R$: $E/L=Gb^2\ln(R/r_0)/[4\pi(1-\nu)]$; the pre-logarithmic
factor is the line tension in the limit of a long, weakly curved
dislocation:
$T=Gb^2/[4\pi(1-\nu)]$. A source segment of length $L$ bowing into
a semicircle of radius $R=L/2$ has arc
length $\pi L/2$:
\begin{equation}
  E_\text{bow}=T\cdot\pi R=\frac{Gb^2L}{8(1-\nu)}.
  \label{eq:Ebow}
\end{equation}

\paragraph{Step~2 (external work)}
Enclosed area $\pi R^2/2=\pi L^2/8$:
\begin{equation}
  W_\text{ext}=\tau_\text{FR}\cdot b\cdot\frac{\pi L^2}{8}.
  \label{eq:Wext}
\end{equation}

\paragraph{Step~3 ($W_\text{ext}=E_\text{bow}$)}
Dividing both sides by $bL/8$:
$\tau_\text{FR}\cdot\pi L=Gb/(1-\nu)$:
\begin{equation}
  \boxed{\tau_\text{FR}=\frac{Gb}{\pi(1-\nu)L}\approx\frac{Gb}{2L}.}
  \label{eq:FR}
\end{equation}
Cu at $L=1\;\mu$m: 14~MPa vs.\ classical 12.4~MPa (13\%).

\subsection{Stacking-fault energy}
\label{sec:SFE}

\subsubsection{FCC Burgers ratios}

A stacking fault in an FCC crystal is bounded by two Shockley
partial dislocations, each carrying only a fraction of the full
lattice Burgers vector; the \OSTZ\ cluster needed to build a partial
is correspondingly smaller than $N_c$. The first step is therefore
to fix that fraction from FCC crystallography alone:
$|b_\text{full}|=a_0/\sqrt2$; $|b_\text{partial}|=a_0/\sqrt6$:
\begin{equation}
  \frac{|b_\text{partial}|}{|b_\text{full}|}
  =\frac{\sqrt2}{\sqrt6}=\frac{1}{\sqrt3}.
  \label{eq:bratio}
\end{equation}

\subsubsection{\OSTZ\ cluster for the Shockley partial}

Setting $N_p\gz W=|b_\text{partial}|=N_c\gz W/\sqrt3$:
$N_p=N_c/\sqrt3$.

Evaluating $(N_c-N_p)/N_p$:
\begin{align*}
  N_c-N_p&=N_c\!\left(1-\frac{1}{\sqrt3}\right),\\
  \frac{N_c-N_p}{N_p}&=\sqrt3\!\left(1-\frac{1}{\sqrt3}\right)=\sqrt3-1.
\end{align*}

\paragraph{Physical picture} When only $N_p$ (rather than the full
$N_c$) \OSTZ s have activated along a glide plane, the crystal has
undergone a partial displacement and the stacking sequence between
the two partial dislocations is faulted
(ABCABC$\to$ABCAB\underline{C}): an intrinsic stacking fault. The
remaining $N_c-N_p$ \OSTZ s have not activated; their unreleased
energy, $(N_c-N_p)\dF$, is stored in the fault.

\subsubsection{SFE derivation}

The fault ribbon has cross-sectional area $\pi W^2$ per \OSTZ\ site
and $N_p$ sites are involved, so the energy per fault area
($\pi W^2$ per OSTZ site, $N_p$ sites):
\[
  \gamma_\text{SF}=\frac{(N_c-N_p)\dF}{\pi W^2 N_p}
  =\frac{(\sqrt3-1)\dF}{\pi W^2}.
\]
Substituting $\dF=\pi\beta_1\gz^2GW^3/3$ and cancelling
$W^3/W^2=W$:
\begin{equation}
  \boxed{\gamma_\text{SF}=\frac{(\sqrt3-1)\beta_1\gz^2GW}{3}.}
  \label{eq:SFE}
\end{equation}
Here $\gz$ is the metal's \emph{ideal shear strain}: taking the
literature value $\gz=0.137$ for Cu~\cite{ogata2004} ($G=48.3$~GPa,
$W=b=0.2556$~nm, $\beta_1=1$) gives $\gamma_\text{SF}=57$~mJ/m$^2$
versus experiment $45$~mJ/m$^2$. Equation~\eqref{eq:SFE} makes
$\gamma_\text{SF}\propto\gz^2Gb$, so the material-to-material spread is
carried by the ideal shear strain, not by a single universal number.
Using the per-metal ideal shear strains of Ogata et~al.~\cite{ogata2004}
\emph{with no fitting}, the predictions for Cu, Al, Ni, Ag, and Au agree
with experiment to within a factor $\sim2$--$3$; the full comparison,
plotted alongside the Peierls stress and core energy of the same metals,
is given in the Supporting Information (Fig.~S2, with the literature
inputs and their sources in Table~S1). The same parameter-free formula
applies to the basal-plane stacking fault of an HCP metal: the SI finds
agreement for Mg but an underprediction for Ti, consistent with the HCP
basal/prismatic anisotropy and the twin-boundary discrepancy of
\S\ref{sec:gb}. No quantitative test is made for BCC metals, which lack a
single well-defined low-energy stacking fault analogous to FCC and HCP.

%%=========================================================
\section{Grain-Boundary Energy}
\label{sec:gb}

A grain boundary is not a single \OSTZ\ object but one of two
distinct elastic configurations, and using the wrong one is a
common source of order-of-magnitude error (\S\ref{sec:RS}).
A \emph{coherent} interface (a twin boundary, or a stacking fault)
stores a coherency strain smeared over its area, so its energy
scales quadratically in the eigenstrain, exactly like the
stacking-fault energy of \S\ref{sec:SFE}. An \emph{incoherent}
low- or high-angle boundary, by contrast, is physically an array of
dislocations, each one an $N_c$-\OSTZ\ chain, so its energy scales
linearly in $b$ through the dislocation line energy. The two
subsections below derive each case from the \OSTZ\ machinery
already established.

\subsection{Coherent boundaries (quadratic in eigenstrain)}

A coherent twin boundary (CTB) carries a coherency strain over an
interfacial area. Its energy is \emph{quadratic} in the eigenstrain,
identical to the \OSET\ SFE~\eqref{eq:SFE}:
$\gamma_\text{CTB}=(\sqrt3-1)\beta_1\gz^2GW/3$.

\subsection{Incoherent boundaries: the Read--Shockley law}
\label{sec:RS}

A low/high-angle boundary is a dislocation array spaced $D=b/\theta$.
In \OSET\ each dislocation is an $N_c$-\OSTZ\ chain; summing the elastic
dipole fields reproduces the classical dislocation line energy
$E_d=(Gb^2/4\pi(1-\nu))\ln(R/r_0)$. Dividing by $D$ with
$R\sim D/2$ ($\ln(R/r_0)\to A-\ln\theta$) yields the
Read--Shockley law~\cite{readshockley1950}:
\begin{align}
  \gamma_\text{GB}(\theta)&=E_0\,\theta(A-\ln\theta),
    \label{eq:RS}\\
  E_0&=\frac{Gb}{4\pi(1-\nu)},\quad A=1+\ln\theta_m.\notag
\end{align}
The GB energy scales linearly with $b$ (dislocation array = linear in
eigenstrain via $b=N_c\gz W$), not quadratically with $\gz$
(activation barrier). Conflating these objects causes 5--40$\times$
error.

Both predictions, $\gamma_\text{CTB}$ and $\gamma_\text{GB}(\theta)$,
were checked for five FCC and two BCC metals in
\S\ref{sec:SFE}/\S\ref{sec:RS} above; the Supporting Information
(Table~S3, Figure~S3) extends this verification to two hexagonal close-packed
metals, Mg and Ti, using the isotropic-elasticity approximation.
The high-angle prediction remains within a factor of $2.2$ for Ti,
consistent with the cubic metals, while the coherent twin-boundary
prediction underpredicts the measured Mg and Ti twin-boundary energy
by a larger margin (up to $6.8\times$), attributable to the
isotropic approximation being less accurate for the strongly
anisotropic, single-plane HCP twinning shear than for cubic-metal
stacking faults.

%%=========================================================
\section{\OSTZ\ Hamiltonian and Statistical Mechanics}
\label{sec:stat}

So far each \OSTZ\ has been treated as an isolated event. In a real
grain boundary or glass, many potential \OSTZ\ sites coexist, each
with its own (disordered) activation barrier, and they interact
elastically through the very dipole field derived in
\S\ref{sec:dipole}. This is the same physical situation as spins in
a disordered magnet: a population of two-state objects (activated or
not) with site-to-site energetic disorder and a long-range
interaction between occupied sites. Casting the \OSTZ\ population in
this language lets us import the standard machinery of mean-field
statistical mechanics to derive the macroscopic strain-rate law and
the strain-hardening behaviour, rather than postulating them
empirically.

\subsection{Construction of the Hamiltonian}
\label{sec:Hamiltonian}

The grain boundary (or, equally, a population of free-volume sites
in a glass) is modelled as a two-dimensional array of discrete
sites, each site $i$ a potential \OSTZ\ location. A binary occupation
number $n_i\in\{0,1\}$ records whether site $i$ is inactive ($n_i=0$)
or activated ($n_i=1$, an \OSTZ\ present there). The total energy of
a configuration $\{n_i\}$ is built additively from three physically
distinct contributions, each derived from results already established
in this paper.

\paragraph{Term~1: Self-energy}
When site $i$ activates, the surrounding elastic matrix must
accommodate the eigenstrain $\gz$ at the cost of the activation
energy $\Delta F_{0,i}$ derived in \S\ref{sec:energy}
(Eq.~\eqref{eq:dF0}). In a real boundary this cost varies from site
to site, since local fluctuations in the free-volume radius $W_i$ and
shear modulus $G_i$ make $\Delta F_{0,i}$ a disordered quantity rather
than a single constant; this disorder is precisely the quenched
randomness later responsible for the lognormal activation-energy
distribution of \S\ref{sec:sinh}. Summing the cost over all sites,
weighted by their occupation,
\begin{equation}
  \mathcal H_\text{self}=\sum_i\Delta F_{0,i}\,n_i.
  \label{eq:Hself}
\end{equation}
(In the mean-field treatment of \S\ref{sec:MF} the site-dependence is
eventually replaced by the disorder-averaged value, $\Delta
F_{0,i}\to\dF$.)

\paragraph{Term~2: Mechanical work}
An applied shear stress $\tau$ performs mechanical work on the
\OSTZ\ as it activates: for a single \OSTZ\ undergoing shear
eigenstrain $\gz$ over volume $V_0$, the work done by the net driving
stress is $W_\mathrm{mech}=(\tau-\tau_0)\gz V_0$, where $\tau_0$ is
the threshold back-stress that must first be overcome to nucleate
the mesoscopic cooperative interface between two neighbouring grains,
$\tau_0\propto L^{-1/2}$~\cite{pad-part1}; $\tau_0$ is therefore
distinct from the elastic interaction stress of Term~3 below. Because
this work \emph{lowers} the energy required to activate the site once
$\tau>\tau_0$, it enters the Hamiltonian with a negative sign and is
summed over all activated sites:
\begin{equation}
  \mathcal H_\text{mech}=-\sum_i(\tau-\tau_0)\gz V_0\,n_i.
  \label{eq:Hmech}
\end{equation}

\paragraph{Term~3: Elastic interaction}
Two activated \OSTZ s at $\mathbf x_i$ and $\mathbf x_j$ also
interact, because the dipole stress field of one (Eq.~\eqref{eq:dipole})
does work on the eigenstrain of the other. Generally, the interaction
energy is the work done by the stress field of \OSTZ\ $i$, evaluated
at the location of \OSTZ\ $j$, against the eigenstrain of \OSTZ\ $j$:
$E_{ij}^\text{int}=-\sigma_{kl}^{(i)}(\mathbf x_j)\,\varepsilon^*_{kl}V_0$.
Since the only nonzero eigenstrain components are
$\varepsilon^*_{13}=\varepsilon^*_{31}=\gz/2$, the contraction picks
out the $13$ and $31$ terms of the stress tensor, which are equal by
symmetry, so the factor of $\half$ from each cancels against the
factor of $2$ from summing the two equal terms, leaving
\begin{equation}
  E_{ij}^\text{int}=-\sigma_{13}^{(i)}(\mathbf x_j)\gz V_0
  =-\frac{G\gz^2V_0^2\mathcal T_{13}(\hat{\mathbf r}_{ij})}
  {2\pi(1-\nu)r_{ij}^3},
  \label{eq:Eint}
\end{equation}
where the second equality substitutes the dipole stress of
Eq.~\eqref{eq:dipole} evaluated at $\mathbf r_{ij}=\mathbf
x_j-\mathbf x_i$. Defining the interaction kernel as
$J_{ij}\equiv-E_{ij}^\text{int}$ (so that a positive $J_{ij}$
corresponds to an energy-\emph{lowering}, cooperative interaction,
matching the usual Ising sign convention) and keeping only the
leading far-field term of $\mathcal T_{13}$,
\begin{align}
  J_{ij}&=G\gz^2V_0^2\,\mathcal K\!\left(\frac{r_{ij}}{W}\right),
    \label{eq:Jij}\\
  \mathcal K(\rho)&=\frac{3\cos^2\theta_{ij}-1}{4\pi(1-\nu)\rho^3}.
    \notag
\end{align}
$J_{ij}>0$ (cooperative) for \OSTZ s aligned along the slip direction
($\theta_{ij}=0$); $J_{ij}<0$ (anti-cooperative) perpendicular.
The form is identical to the magnetic dipole--dipole interaction: a
direct consequence of the shared $r^{-3}$ Green's function structure.
Summing this pairwise interaction over all activated pairs, with the
conventional factor $\half$ inserted to avoid double-counting each
pair $(i,j)$ and $(j,i)$ separately,
\begin{equation}
  \mathcal H_\text{int}=-\half\sum_{i\ne j}J_{ij}\,n_in_j.
  \label{eq:Hint}
\end{equation}

\paragraph{Full Hamiltonian}
Adding the three terms gives the total energy of an arbitrary
occupation configuration $\{n_i\}$:
\begin{multline}
  \mathcal H=\sum_i\!\left[\Delta F_{0,i}-(\tau-\tau_0)\gz V_0\right]n_i\\
  -\half\sum_{i\ne j}J_{ij}\,n_in_j.
  \label{eq:Hamiltonian}
\end{multline}
This is formally identical to a \emph{random-field Ising model} in a
uniform external field $h=(\tau-\tau_0)\gz V_0$, with quenched random
self-energies $\Delta F_{0,i}$ playing the role of site disorder and
the cooperative coupling $J_{ij}$ playing the role of the
ferromagnetic exchange interaction. This identification is what
allows the standard mean-field machinery of magnetic statistical
mechanics to be imported wholesale into the \OSTZ\ problem in the
next subsection.

\subsection{Mean-field decoupling}
\label{sec:MF}

\paragraph{The linearisation step}
Write $n_i=\bar n+\delta n_i$. Expand $n_in_j$ and drop
$\delta n_i\delta n_j$ (Bragg--Williams approximation~\cite{bragg1934}):
\begin{equation}
  n_in_j\approx n_i\bar n+n_j\bar n-\bar n^2.
  \label{eq:BW}
\end{equation}

\paragraph{Effect on the interaction term}
Substituting Eq.~\eqref{eq:BW} into the interaction term
$-\half\sum_{i\ne j}J_{ij}n_in_j$: the first two
terms are equal by relabelling ($i\leftrightarrow j$), and with
translation invariance $\sum_{j\ne i}J_{ij}=zJ_0$:
\begin{equation}
  -\half\sum_{i\ne j}J_{ij}n_in_j\approx
  -z\bar nJ_0\sum_i n_i+\text{const.}
  \label{eq:MFinteraction}
\end{equation}

\paragraph{Mean-field Hamiltonian}
The total factorises into independent single-site problems:
\begin{align}
  \mathcal H^\text{MF}&=\sum_i h^\text{eff}n_i+\text{const},
    \label{eq:heff}\\
  h^\text{eff}&=\dF-(\tau-\tau_0)\gz V_0-z\bar nJ_0.\notag
\end{align}
The effective barrier $h^\text{eff}$ is the bare cost $\dF$ reduced
by the mechanical driving force and the cooperative elastic
interaction from already-activated neighbours.

\subsubsection{Self-consistency equation}
\label{sec:selfconsistency}

Single-site partition function $\mathcal Z_i=1+e^{-h^\text{eff}/kT}$.
Mean occupation:
\begin{equation}
  \bar n=\langle n_i\rangle
  =\frac{e^{-h^\text{eff}/kT}}{1+e^{-h^\text{eff}/kT}}.
  \label{eq:selfconsistency}
\end{equation}
Substituting $h^\text{eff}$ explicitly:
\begin{equation}
  \boxed{\bar n=\frac{\exp\!\left[-\tfrac{\dF-(\tau-\tau_0)\gz V_0
    -z\bar nJ_0}{kT}\right]}{1+\exp\!\left[-\tfrac{\dF-(\tau-\tau_0)
    \gz V_0-z\bar nJ_0}{kT}\right]}}
  \label{eq:selfconsfull}
\end{equation}
This nonlinear equation is solved iteratively for $\bar n$.

In the \emph{dilute limit} ($h^\text{eff}\gg kT$, $z\bar nJ_0\to0$):
\[
  \bar n\approx\exp(-\dF/kT)\cdot\exp\!\left(\frac{(\tau-\tau_0)\gz V_0}{kT}\right),
\]
reflecting only the forward activation. The full forward--backward balance
is captured by the Eyring kinetic route (\S\ref{sec:sinh}).

\subsection{Sinh flow law: step-by-step derivation}
\label{sec:sinh}

\paragraph{Step~1 (strain per activation)}
Each \OSTZ\ activation shifts the material by $\beff=\gz W$ over a
cross-section $\pi W^2$ in a grain of size $d$. The macroscopic
strain increment:
\begin{equation}
  \delta\gamma_{0,\text{macro}}=\frac{\beff\cdot\pi W^2}{\pi W^2\cdot d}
  =\frac{\gz W}{d}.
  \label{eq:deltaGamma}
\end{equation}

\paragraph{Step~2 (Eyring forward/backward rates)}
Following Eyring~\cite{eyring1935}, bias $\Delta w=(\tau-\tau_0)\gz V_0/2$.
Forward barrier $\dF-\Delta w$; reverse $\dF+\Delta w$:
\begin{equation}
  \Gamma^\pm=\nu_0\exp\!\left(-\frac{\dF\mp\Delta w}{kT}\right).
  \label{eq:Gamma}
\end{equation}

\paragraph{Step~3 (factor out $e^{-\dF/kT}$)}
Using $e^{a+b}=e^a\cdot e^b$:
\begin{align*}
  e^{-(\dF-\Delta w)/kT}&=e^{-\dF/kT}\cdot e^{+\Delta w/kT},\\
  e^{-(\dF+\Delta w)/kT}&=e^{-\dF/kT}\cdot e^{-\Delta w/kT}.
\end{align*}

\paragraph{Step~4 ($e^x-e^{-x}=2\sinh x$)}
\begin{multline}
  \dot N_\text{net}=\Gamma^+-\Gamma^-\\
  =2\nu_0\sinh\!\left[\frac{(\tau-\tau_0)\gz V_0}{2kT}\right]
  e^{-\dF/kT}.
  \label{eq:Nnet}
\end{multline}

\paragraph{Step~5 (macroscopic rate)}
Multiply by $\delta\gamma_{0,\text{macro}}=\gz W/d$:
\begin{equation}
  \boxed{\dot\gz=\frac{2W\gz\nu_0}{d}
  \sinh\!\left[\frac{(\tau-\tau_0)\gz V_0}{2kT}\right]
  e^{-\dF/kT}.}
  \label{eq:rate}
\end{equation}
Identical to Padmanabhan et al.\ (1996, Eq.~8)~\cite{pad-part1}.

\subsubsection{Comparison: partition-function route}

From $\ln\mathcal Z=N_\text{sites}\ln(1+e^{-h^\text{eff}/kT})$,
differentiation with respect to $(\tau-\tau_0)\gz V_0/kT$ gives
$N_\text{sites}\bar n$. In the dilute limit, this is a
\emph{forward-bias-only} expression, not equal to $2e^{-\dF/kT}
\sinh(\Delta w/kT)$.
The sinh rate equation is an Eyring \emph{kinetic} result.
Detailed balance confirms this: $\Gamma^+/\Gamma^-=e^{2\Delta w/kT}$.
Figure~S4 in the Supporting Information compares the correct
partition-function occupation against the naive $\sinh$ form
sometimes used in the literature, quantifying the error of the
latter.

\subsubsection{Connection to Padmanabhan et al.\ (1996)}

With $\Delta w=(\tau-\tau_0)\gz V_0/2$ and $(2W\gz/d)\dot w=\dot\gz$,
Eq.~(8) of~\cite{pad-part1} is recovered exactly.
Including a lognormal disorder distribution $f(\dF)$:
\begin{multline}
  \dot\gz=\frac{2W\gz\nu_0}{d}\int_0^\infty f(\dF)\\
  \times\sinh\!\left[\frac{(\tau-\tau_0)\gz V_0}{2kT}\right]
  e^{-\dF/kT}\,\mathrm d(\dF).
  \label{eq:rateDisorder}
\end{multline}
Universal lognormal parameter $\bar A=5.6021$~\cite{venkatesh-part4};
CV$=14.8\%$; range $[0.28,0.52]$~eV.

\subsection{\OSTZ\ interaction as the origin of Taylor hardening}
\label{sec:taylor}

When $z\bar nJ_0$ is non-negligible, the effective threshold
stress increases:
\begin{equation}
  \tau_0^\text{eff}=\tau_0+\frac{z\bar nJ_0}{\gz V_0}.
  \label{eq:tau0eff}
\end{equation}
With $\bar n=\pi W^2\rho_\text{disl}$ and
$J_0\approx G\gz^2V_0^2/[4\pi(1-\nu)W^3]$,
substituting $V_0=(2\pi/3)W^3$:
\begin{equation}
  \tau_0^\text{eff}-\tau_0=
  \frac{z\gz G\pi W^2}{6(1-\nu)}\,\rho_\text{disl}.
  \label{eq:Taylorhardening}
\end{equation}
\OSET\ predicts hardening \emph{linear} in $\rho_\text{disl}$
(stage-I regime). The classical Taylor $\sqrt\rho$ law~\cite{taylor1934} requires
the additional geometric forest-obstacle argument
($\tau_c=Gb/L$, $L\propto\rho^{-1/2}$) and is not derivable
from the \OSTZ\ Ising Hamiltonian alone (Figure~S5 in the Supporting
Information compares the two hardening laws directly).

%%=========================================================
\section{Why \OSET\ is More Fundamental}
\label{sec:fundamental}

\subsection{Logical priority theorem}

The results assembled in this paper can be stated as a precise
logical claim about the relationship between the two theories,
analogous to the way a more general physical theory contains a more
restricted one as a special limit (as, for instance, relativistic
mechanics contains Newtonian mechanics in the limit $v\ll c$).

\paragraph{Theorem} \OSET\ logically entails classical dislocation
theory, but classical dislocation theory does not entail \OSET; that
is, \OSET$\Rightarrow$dislocation theory holds, while the converse
implication dislocation theory$\Rightarrow$\OSET\ fails.

\paragraph{Proof of $\Rightarrow$} Sections~\ref{sec:PN}--\ref{sec:chainStress}
show, by direct construction and without assuming the existence of a
dislocation anywhere in the argument, that a co-planar chain of $N$
\OSTZ s reproduces every defining property of a dislocation: the
exact \PN\ displacement profile (\S\ref{sec:PN}), a non-singular
stress field that reduces to the singular Volterra field
$\sigma\sim Gb/2\pi r$ in the formal limit $N\to\infty$, $W\to0$ at
fixed $b=N\beff$ (\S\ref{sec:chainStress}), and a Burgers circuit
that closes on a quantised lattice vector $b=N_c\gz W\in
\Lambda_\mathrm{Bravais}$ precisely when the chain reaches the
critical length $N_c$ derived in \S\ref{sec:Nc}. Since every
quantity that the Burgers-circuit definition of
Eq.~\eqref{eq:burgers} requires (a singular line, a quantised
Burgers vector, an athermal core) is obtained as a derived,
large-$N$ limiting property of the \OSTZ\ chain, the classical
dislocation is recovered as a special case of \OSET\ rather than
assumed as an independent postulate. $\square$

\paragraph{Proof that the converse fails} It suffices to exhibit
physical situations that \OSET\ describes but that the
Burgers-circuit definition cannot, since by definition a theory
cannot entail phenomena that lie entirely outside its domain of
applicability. Four such situations were identified in
\S\ref{sec:intro}--\S\ref{sec:classes}: (a)~a single, sub-critical
\OSTZ\ event ($N<N_c$), which carries a well-defined activation
energy and stress dipole (\S\ref{sec:dipole}) but no Burgers vector,
since no closed circuit around it has a non-zero closure failure;
(b)~thermally activated shear in amorphous matter, where the
absence of any Bravais lattice $\Lambda_\mathrm{Bravais}$ in
Eq.~\eqref{eq:burgers} makes the dislocation construction
inapplicable by definition, not merely inconvenient
(\S\ref{sec:classes}); (c)~grain-boundary sliding at a general
high-angle boundary, where the DSC lattice required to define a
boundary dislocation is commensurate with the two adjoining crystals
only statistically (\S\ref{sec:intro}); and (d)~the continuous
nanocrystal-to-glass crossover as $d\to d_0$, where the effective
cooperative number $N_c^\text{eff}(d)\to1$ so that the very notion of
a multi-\OSTZ\ chain (and hence of a dislocation) ceases to be
meaningful, while the underlying \OSTZ\ description remains valid
throughout (\S\ref{sec:classes}). In each case \OSET\ supplies a
well-posed description with no discontinuity in its governing
equations, while dislocation theory supplies none. $\square$

\subsection{Ontological hierarchy}

\begin{center}
\small
$\text{OSTZ}
\xrightarrow{N=1}\text{STZ/GBS}
\xrightarrow{N=N_p}\text{partial disln}$\\[3pt]
$\xrightarrow{N=N_c}\text{full disln}
\xrightarrow{N\gg N_c}\text{slip band}$
\end{center}

Dislocations occupy one branch of the \OSTZ\ tree. \OSET\ is the root.

\begin{table}[htbp]
\centering
\small
\caption{Classical dislocation theory vs.\ \OSET.}
\label{tab:comparison}
\renewcommand{\arraystretch}{1.35}
\begin{tabular}{p{1.5cm}p{2.2cm}p{2.3cm}}
\toprule
Property & Classical & \OSET\\
\midrule
Entity & Topological line & Eshelby spheroid\\
Lattice & Required & Not required\\
Materials & Crystals & All\\
Core stress & $\to\infty$ & $\le G\gz$\\
Core energy & Fitted $\alpha Gb^2$ & Eq.~\eqref{eq:coreE}\\
Core width & Fitted & $\zeta=W$\\
Burgers $b$ & Input & $N_c\gz W$\\
Thermal act. & Appended & Intrinsic\\
$\tau_P$ & $W$ fitted & Eq.~\eqref{eq:peierls}\\
$\tau^*$ & Frenkel sep. & Eq.~\eqref{eq:strength}\\
SFE & Fitted & Eq.~\eqref{eq:SFE}\\
GB plast. & DSC needed & Natural\\
Quantum repn. & Topological field theory & Point-particle Schr.\ eq.\ (\S\ref{sec:quantum})\\
\bottomrule
\end{tabular}
\renewcommand{\arraystretch}{1}
\end{table}

\subsection{A quantum-mechanical representability argument}
\label{sec:quantum}

The theorem above is a purely logical statement: \OSTZ\ generates
dislocation theory as a special case, and not conversely. A second,
independent argument for the same hierarchy follows from asking
which of the two objects can, in principle, be generated by an
ordinary (first-quantised) Schr\"odinger equation for a particle
moving on a finite-dimensional configuration space, exactly the
description routinely used for a point defect such as a vacancy or
self-interstitial.

With $W$ set to the free-volume length scale, and in particular in
the limit $W\to b$ already adopted for the stacking-fault and
core-energy predictions of \S\ref{sec:SFE}--\S\ref{sec:coreE}, the
\OSTZ\ of Eq.~\eqref{eq:geometry} shrinks to an atomic-scale,
compact inclusion: a finite-volume eigenstrained region located at a
single site, in exactly the same sense that a vacancy or
interstitial is treated, in continuum defect theory, as a point
centre of dilatation despite its finite atomic core. This is not an
additional assumption: it is the same idealisation already used,
without further approximation, throughout \S\ref{sec:stat}, where
every \OSTZ\ is represented by a binary occupation variable
$n_i\in\{0,1\}$ at a discrete site $\mathbf x_i$ in
Hamiltonian~\eqref{eq:Hamiltonian}: crystallography's vacancy
bookkeeping applied to a shear excitation. Three properties of this
point-like limit are what make a Schr\"odinger description of its
activation kinetics well posed.

(i) \emph{Finite, system-size-independent self-energy.} The \OSTZ\
stress and displacement fields are those of a finite elastic dipole
(Eq.~\eqref{eq:dipole}), decaying as $r^{-3}$ in the far field
(Eq.~\eqref{eq:Jij}); the resulting activation energy $\dF$
(Eq.~\eqref{eq:dF0}) is a single finite number, fixed entirely by
local quantities ($G$, $\gz$, $\varepsilon_0$, $V_0$) and independent
of the size $R$ of the surrounding crystal.

(ii) \emph{Finite-dimensional configuration space.} Its state is
specified by a finite number of coordinates (occupied or
unoccupied at a given site, or, in a continuum refinement, a small
set of local reaction coordinates along the activation path), not
by a field defined along an extended continuum of points.

(iii) \emph{No topological charge.} A sub-critical \OSTZ\ event
($N<N_c$, proof of the converse above) is simply absent, with no
analogue of a circuit integral required to detect it; "activated" or
"not" is a local, not a non-local, property.

These three properties are precisely the conditions under which a
localised excitation can be assigned a genuine potential-energy
surface $U(\mathbf q)$, finite and independent of system size,
over a finite-dimensional configuration space $\mathbf q$, so
that its quantum dynamics is generated by an ordinary Schr\"odinger
equation of exactly the kind used throughout solid-state defect
physics to treat vacancy formation and migration
(configuration-coordinate and multiphonon theory). The classical
Eyring/Kramers rate law already used for the \OSTZ,
Eq.~\eqref{eq:rate} of \S\ref{sec:sinh}, is recovered as the
high-temperature limit $\hbar\omega\ll kT$ of exactly this problem,
so the \OSTZ\ kinetic theory developed in this paper rests on the
same quantum-statistical foundation as point-defect diffusion
theory, even though only its classical limit is used here.

A dislocation line fails all three conditions simultaneously.
(i)~Its self-energy diverges logarithmically with the system size,
$E_d=Gb^2/4\pi(1-\nu)\ln(R/r_0)$ (\S\ref{sec:RS}): no
system-size-independent potential $U$ can be assigned to the line by
itself, since its energy depends explicitly on the extent of the
surrounding crystal and not only on the defect's own coordinates.
(ii)~Its configuration is a continuous field $\mathbf u(s)$ defined
along the entire arc length $s$ of the line: an
infinite-dimensional configuration space, not the finite set of
coordinates a point-like object admits. (iii)~Its defining
invariant is the topological closure failure of a circuit integral,
Eq.~\eqref{eq:burgers}, a non-local property of the displacement
field around the whole line, with no representation as the local
amplitude of a wavefunction at a point. Each obstruction is,
individually, sufficient to make a single-particle Schr\"odinger
description ill posed: a divergent, system-size-dependent
self-energy precludes a finite potential well; an
infinite-dimensional configuration space precludes finite-dimensional
quantum mechanics; and a non-local topological charge has no local
wavefunction representation. A quantum treatment of a dislocation is
not impossible in principle, but it requires the qualitatively
different machinery of a quantised field with topological
(winding-number) excitations, the same formalism used for vortex
lines in superfluids or flux lines in superconductors, rather
than the elementary Schr\"odinger equation that suffices for the
point-like \OSTZ.

This furnishes a third, physical argument for the hierarchy
established above: only the point-like \OSTZ, not the line-like
dislocation built from it, reduces to an ordinary
quantum-mechanical problem of the kind already standard for vacancies
and other point defects. The dislocation inherits a well-defined
activation kinetics only because it is, by construction
(\S\ref{sec:chainStress}--\S\ref{sec:Nc}), itself a finite chain of
such point-like objects; considered in isolation, as an
infinitely extended elastic line, it has no analogous
single-particle quantum-mechanical foundation.

%%=========================================================
\section{\OSET\ Across Material Classes}
\label{sec:classes}

The single \OSTZ\ framework developed above behaves differently
depending on how easily \OSTZ s can cluster, and this one parameter
is enough to span the full range of material classes encountered in
plasticity. In a perfect crystal, where $\dF\gg kT$, individual
\OSTZ\ activations are rare and the lattice forces them to act in
cooperative clusters of exactly $N_c$, so the classical Orowan
equation~\cite{orowan1940} is recovered as the large-cluster limit
of the unified rate law. At a grain boundary or in a superplastic alloy, by contrast,
\OSTZ s nucleate at free-volume sites with $N\ll N_c$, and the
resulting rate equation reduces directly to Eq.~(39) of
Padmanabhan et al.~\cite{pad-part1}. In a metallic glass there is no
periodic lattice to enforce any particular cluster size at all, so
$N_c$ is undefined and deformation proceeds through individual
\OSTZ\ activations, exactly the shear-transformation-zone
picture~\cite{argon1979,falk-langer1998}: at low temperature this
produces localised deformation, with shear bands forming as
percolating paths of activated \OSTZ s, while at high temperature
the activations are dense enough that flow becomes homogeneous.
Finally, as a nanocrystal's grain size $d$ shrinks toward
$d_0=2\sqrt6\,W$, the threshold stress $\tau_0\to0$ and the
effective cooperative number $N_c^\text{eff}(d)\to1$, so the
crystal smoothly crosses over into glass-like behaviour. \OSET\
therefore spans the nanocrystal-to-glass transition naturally,
without switching to a different model at the crossover.

%%=========================================================
\section{Unified Rate Equation}
\label{sec:unified}

The sinh flow law of \S\ref{sec:sinh} was derived for a single,
sharp activation barrier and a dilute population of sites; real
materials have a distribution of barriers (captured by $\sigma_F^2$)
and, in crystals, a finite cooperative number $N_c$ rather than the
GB limit $N\to1$. The equation below restores both effects,
combining the \OSTZ\ partition function (\S\ref{sec:stat}) with
disorder-field theory (activation energy variance $\sigma_F^2$)
and crystalline crossover $\Theta$:
\begin{multline}
  \dot\gz=\frac{2W\bar\gz\nu_0}{d}
  \sinh\!\left[\frac{(\tau-\tau_0^*)\bar\gz V_0}{2kT}\right]\\
  \times\exp\!\!\left(-\frac{\Delta\bar F_0}{kT}
  +\frac{\sigma_F^2}{2(kT)^2}\right)\Theta(N_c,T,d),
  \label{eq:rateUnified}
\end{multline}
where $\tau_0^*=\tau_0\sqrt{1+\rho_{G\gz}\sigma_G\sigma_{\gz}}$ and
\begin{equation}
  \Theta=\begin{cases}
    1 & \text{glass/GB regime}\\
    N_c^{-1}\!\sum_{N=1}^{N_c}\!N\,e^{-E_c(N)/kT} & \text{crystal}
  \end{cases}
  \label{eq:Theta}
\end{equation}
In the crystal limit $\Theta\to N_c$ (Orowan); in the glass/GB limit
$\Theta\to1$ (Padmanabhan et al.~\cite{pad-part1}, Eq.~(39)).

%%=========================================================
\section{Validation}
\label{sec:validation}

\OSET\ is validated below at three increasingly independent levels:
against the original 1996 Padmanabhan et al.\ parameter set, against
literature published independently after 1996, and against the
recent 41-system compilation of Harisankar and
Padmanabhan~\cite{harisankar2025}. Full numerical tables underlying
every comparison summarised here are given in the Supporting
Information.

The Eshelby-tensor calculation of \S\ref{sec:eshelby} reproduces the
constraint factors of the original theory exactly: $\beta_1=0.446$
and $\beta_2=0.889$. Using these canonical values, the activation
energy predicted from Eq.~\eqref{eq:dF0} is $\dF=0.405$~eV, against
the original $0.38$~eV, agreement to within 6\%. The dilatational
eigenstrain $\varepsilon_0$ is likewise consistent with the
underlying free-volume picture: for $W=2.5b$, the predicted free
volume per activation, $\delta V_\text{free}=\varepsilon_0
V_0\approx0.1\,\Omega$ per grain-boundary atom, confirms the
molecular-dynamics estimate of Wolf~\cite{wolf1990}. The critical
cluster number is reproduced exactly as well: $N_c=b/(\gz
W)=0.3/(0.1\times0.75)=4.0$, matching the cooperative cluster size
of $N=4$ measured experimentally in Zn--22Al~\cite{astanin-part2}.
The full rate equation, evaluated for Al-12Si at $773$~K with
$d=10\;\mu$m and $\tau-\tau_0=5$~MPa, predicts
$\dot\gz=0.16$~s$^{-1}$, just above the upper edge of the typical
superplastic strain-rate window $10^{-3}$--$10^{-1}$~s$^{-1}$; the
threshold-stress scaling $\tau_0\propto L^{-1/2}$ is confirmed for
Al-33Cu to within $10\%$; and the lognormal disorder distribution
introduced in \S\ref{sec:sinh} reproduces the universal fitted
parameters $\bar A=5.6021$ and $\mathrm{CV}=14.8\%$ (Table~S4 in the
Supporting Information gives the complete claim-by-claim breakdown).

Independent literature published after the original 1996 papers
provides a broader test of the same rate equation and parameter
set. Sripathi \& Padmanabhan~\cite{sripathi2014} applied it to 33
distinct material systems and obtained correlation coefficients of
$R=0.89$--$1.00$ throughout, and in doing so introduced an
alternative $\beta_1$ convention that is fully reconciled with the
one used here once the corresponding $\gz$ is rescaled. The canonical
value $\gz=0.10$ is itself supported by four independent methods,
spanning molecular dynamics~\cite{wolf1990}, bubble-raft analogue
experiments, shear-band initiation in metallic
glasses~\cite{argon1979}, and shear-transformation-zone
theory~\cite{falk-langer1998}. The same framework extends to bulk
metallic glasses~\cite{buenz2015}, nanocrystalline
nickel~\cite{pad2018}, and 8 further high-strain-rate superplastic
alloy and composite systems~\cite{padbasariya2009}, in every case to
within a factor of $2$--$2.5$. The Peierls stress formula of
\S\ref{sec:peierls} is exact for Fe and within a factor of 2 for Cu
when compared against the compilation of
Nabarro~\cite{nabarro1997}, and the Lorentzian dislocation-density
profile of \S\ref{sec:lorentzian} matches X-ray diffraction
line-profile data for Al~\cite{vermeulen1995} (Table~S5 in the
Supporting Information gives the complete comparison, and
Figure~S6 there consolidates the bulk-metallic-glass,
nanocrystalline-Ni, and high-strain-rate-alloy comparisons into a
single predicted-to-measured parity plot alongside the Al-12Si worked
example above).

\label{sec:harisankar}
Table~S6 of the Supporting Information compares \OSET\ against the
mean of the semi-empirical constants fitted by Harisankar and
Padmanabhan~\cite{harisankar2025} across all 41 systems in their
compilation, with the full system-by-system breakdown given there in
Table~S7 and the system-by-system distributions plotted in
Figure~S7. Four findings emerge. The
dilatational eigenstrain is predicted to better than $2\%$, the
tightest agreement of any parameter, confirming that $\varepsilon_0$
is a genuine material constant rather than a fitted quantity. The
mean fitted shear eigenstrain, $\gz=0.085$, lies between the
grain-boundary ($0.10$) and crystal ($0.12$) canonical values used
throughout this paper, consistent with the temperature-softening
trend $\gz(T)=0.0827+1.342/T$ extracted independently in the same
dataset (Figure~S8). The activation energy $\dF$ accounts for only about a
quarter of the total measured activation energy $Q$
($\dF\approx Q/4$); the remaining three-quarters is plausibly
non-elastic in origin, comprising grain-boundary migration,
triple-junction accommodation, and the cooperative rearrangement of
roughly $N_a\approx16$ boundary units. The grain-boundary energy
$\gamma_B$ is reproduced to a median factor of $4.5\times$, and
critically this offset is a constant multiplicative factor showing
no correlation with homologous temperature ($r=0.00$), indicating
that \OSET\ captures the correct $Gb$ scaling even where the
absolute magnitude is offset.

%%=========================================================
\section{Discussion}
\label{sec:discussion}

\OSET\ establishes the ontological hierarchy: single \OSTZ\ = elastic
dipole = GBS/STZ event; $N$-chain = \PN\ dislocation;
$N=N_c$ = lattice dislocation.

\paragraph{Strengths}
\OSET\ performs best for the FCC noble and transition metals (Cu,
Ni, Au, Ag), where the stacking-fault energy (within $5\%$ for Cu),
the Peierls stress (within a factor of 2), and the core energy are
all reproduced without any free parameters, and the grain-boundary
and superplastic rate equation has been validated across 41
material systems.

\paragraph{Limitations}
Three discrepancies are worth flagging openly. Aluminium has an
anomalously high stacking-fault energy that can only be matched by
inflating the eigenstrain to $\gz\approx0.30$, a symptom of its
nearly-free-electron band structure rather than a failure of the
elastic framework. In BCC metals the Peierls stress is
underpredicted by a factor of $10$--$100\times$~\cite{nabarro1997},
because the true BCC dislocation core is narrower than the lattice
spacing ($W<b$); using $W=0.2b$ instead of $W\approx b$ recovers the
experimental value. Finally, the activation energy ratio
$Q/\dF\approx4$ found in \S\ref{sec:harisankar} is consistent with a
mesoscopic cooperative number $N_a\approx16$ ($Q=N_a\dF$) combined
with the lognormal tail of the disorder distribution, indicating
that the purely elastic \OSTZ\ barrier accounts for only about 25\%
of the total measured activation energy.

\paragraph{Outlook}
Two structural predictions of \OSET\ have not yet been tested
directly against microscopy or atomistic simulation: the bounded
core stress $\sigma_\text{core}\le G\gz$ and the core-width
identification $\zeta=W$ (Table~\ref{tab:comparison}). High-resolution
transmission electron microscopy or atomistic simulation of an
isolated dislocation core in a well-characterised FCC metal could
measure both quantities independently of the bulk superplastic-flow
data used in \S\ref{sec:validation}, and so would constitute a
sharper, core-level test of the \OSTZ\ picture. A second open
direction is the extension from the single eigenstrain component
$\varepsilon^*_{13}$ treated throughout this paper to general
multiaxial loading, which would allow \OSET\ to be compared directly
against polycrystal plasticity models rather than only against
uniaxial superplastic-flow data.

\paragraph{Broader implications}
Because \OSET\ requires no periodic lattice, the same three
parameters ($W$, $\gz$, $\varepsilon_0$) describe a crystalline
dislocation core, a metallic-glass shear-transformation zone, and a
grain-boundary sliding event alike. Design strategies already
developed for dislocation-mediated plasticity, such as core-width
control through solute segregation, may therefore transfer in
modified form to amorphous and nanocrystalline systems, where no
comparably quantitative design rule currently exists.

%%=========================================================
\section{Conclusions}
\label{sec:conclusions}

\begin{enumerate}

\item \OSTZ\ is lattice-free, non-singular, thermally activated.
  Eshelby parameters $\beta_1=0.446$, $\beta_2=0.889$ are derived
  rigorously from $I_1$, $I_{13}$
  (Sections~\ref{sec:I1}--\ref{sec:numeval}).

\item Dislocations emerge as \OSTZ\ collectives:
  dipole Eq.~\eqref{eq:dipole}, \PN\ profile Eq.~\eqref{eq:PNprofile},
  non-singular chain stress Eq.~\eqref{eq:chainStress}; Volterra
  singularity only as $W\to0$, Eq.~\eqref{eq:volterra}.

\item Five dislocation observables are derived without adjustable
  parameters: $\tau^*$~\eqref{eq:strength}, $\tau_P$~\eqref{eq:peierls},
  $E_\text{core}$~\eqref{eq:coreE}, $\tau_\text{FR}$~\eqref{eq:FR},
  $\gamma_\text{SF}$~\eqref{eq:SFE}.

\item The \OSTZ\ Hamiltonian~\eqref{eq:Hamiltonian} yields the sinh
  flow law~\eqref{eq:rate} and linear
  hardening~\eqref{eq:Taylorhardening}.

\item Validation: $\varepsilon_0$ within $2\%$, $\gz$ within $15\%$,
  $\gamma_B$ to median $4.5\times$, across 41 superplastic systems.

\item \OSET$\Rightarrow$dislocation theory (not conversely):
  \OSET\ is the more fundamental description
  (Table~\ref{tab:comparison}), a conclusion reinforced by a
  quantum-mechanical representability argument: only the point-like
  \OSTZ, not the topologically extended dislocation line, reduces to
  an ordinary Schr\"odinger-equation problem (\S\ref{sec:quantum}).

\end{enumerate}

\section*{CRediT}
\textbf{A.~Linda:} Methodology, Formal analysis, Writing.
\textbf{K.A.~Padmanabhan:} Conceptualization, Supervision, Review.

\section*{Competing interests}
None.

\section*{Data availability}
A Python notebook with all derivations and numerical verifications
accompanies this article. Every symbolic derivation step in that
notebook (Eshelby tensor integrals, the dipole and glide-plane stress
fields, the stacking-fault and Read--Shockley energy derivations, the
\OSTZ\ Hamiltonian and mean-field rate equation, and others) is
checked independently using the symbolic mathematics library
SymPy~\cite{meurer2017sympy}, alongside NumPy/SciPy for the numerical
cross-checks.

\bibliographystyle{elsarticle-num}
\bibliography{OSET_references}

\end{document}
